% Auto-generated by gen_appendix_tests.py
% Source: metamath-test/run-testsuite-all

\subsection{Character Set (L73--79)}

\begin{quote}\small\itshape
file are the 94 non-whitespace printable ascii characters, which are digits,
upper and lower case letters, and the following 32 special characters:
! " \# \$ \% \& ’ ( ) * + , - . / :
; < = > ? @ [ \textbackslash{} ] \^{} \_ ‘ \{ | \} \~{}
plus the following characters which are the “white space” characters: space
(a printable character), tab, carriage return, line feed, and form feed. We
will use typewriter font to display the printable characters.
\end{quote}

\noindent\textbf{test01\_non-printable\_ascii\_characters.mm} [REJECT] \textit{L73-79: only printable ASCII  +  whitespace allowed}\par

\begin{lstlisting}[language=Metamath,basicstyle=\ttfamily\scriptsize,frame=single,framerule=0.3pt,rulecolor=\color{black!20},backgroundcolor=\color{black!2},xleftmargin=1em,framexleftmargin=0.5em,numbers=none,aboveskip=0.5em,belowskip=0.3em]
$( Unit Test 1: Non-printable ASCII characters $)
$( Should reject: True $)

$c wff |- $.
$c → $.
$v x $.
wf $f wff x $.
ax $a |- x $.
\end{lstlisting}

\subsection{Whitespace Characters (L77--79)}

\begin{quote}\small\itshape
plus the following characters which are the “white space” characters: space
(a printable character), tab, carriage return, line feed, and form feed. We
will use typewriter font to display the printable characters.
\end{quote}

\noindent\textbf{emptyline.mm} [ACCEPT] \textit{L77-79: whitespace (including newlines) is valid separator}\par

\begin{lstlisting}[language=Metamath,basicstyle=\ttfamily\scriptsize,frame=single,framerule=0.3pt,rulecolor=\color{black!20},backgroundcolor=\color{black!2},xleftmargin=1em,framexleftmargin=0.5em,numbers=none,aboveskip=0.5em,belowskip=0.3em]

\end{lstlisting}

\noindent\textbf{test54\_whitespace\_variety.mm} [ACCEPT] \textit{L77-79: various whitespace chars valid}\par

\begin{lstlisting}[language=Metamath,basicstyle=\ttfamily\scriptsize,frame=single,framerule=0.3pt,rulecolor=\color{black!20},backgroundcolor=\color{black!2},xleftmargin=1em,framexleftmargin=0.5em,numbers=none,aboveskip=0.5em,belowskip=0.3em]
$( Test 54: Whitespace variety - tabs and spaces are valid whitespace per Sec 4.1.1; should ACCEPT $)

$c	wff	|- $.
$v	ph	$.

f1	$f	wff	ph	$.
a1	$a	|- ph	$.
p1	$p	|- ph	$=	f1	 a1	$.
\end{lstlisting}

\noindent\textbf{test68\_formfeed\_whitespace.mm} [ACCEPT] \textit{L77-79: form-feed (0x0C) is whitespace}\par

\begin{lstlisting}[language=Metamath,basicstyle=\ttfamily\scriptsize,frame=single,framerule=0.3pt,rulecolor=\color{black!20},backgroundcolor=\color{black!2},xleftmargin=1em,framexleftmargin=0.5em,numbers=none,aboveskip=0.5em,belowskip=0.3em]
$( Test 68: Form-feed (0x0C) is valid whitespace per Metamath spec section 4.1.1; should ACCEPT $)

$cwff|-$.
$vph$.

f1$fwffph$.
a1$a|-ph$.
p1$p|-ph$=f1a1$.
\end{lstlisting}

\subsection{Token Separation (L81--83)}

\begin{quote}\small\itshape
separated by white space (which is any sequence of one or more white
space characters). The set of keyword tokens is \$\{, \$\}, \$c, \$v, \$f, \$e, \$d,
\$a, \$p, \$., \$=, \$(, \$), \$[, and \$]. The last four are called auxiliary or
\end{quote}

\noindent\textbf{comment1.mm} [REJECT] \textit{L81-83: whitespace required between tokens}\par

\begin{lstlisting}[language=Metamath,basicstyle=\ttfamily\scriptsize,frame=single,framerule=0.3pt,rulecolor=\color{black!20},backgroundcolor=\color{black!2},xleftmargin=1em,framexleftmargin=0.5em,numbers=none,aboveskip=0.5em,belowskip=0.3em]
$( $( $)
\end{lstlisting}

\noindent\textbf{test02\_missing\_whitespace\_between\_tokens.mm} [REJECT] \textit{L81-83: tokens must be separated by whitespace}\par

\begin{lstlisting}[language=Metamath,basicstyle=\ttfamily\scriptsize,frame=single,framerule=0.3pt,rulecolor=\color{black!20},backgroundcolor=\color{black!2},xleftmargin=1em,framexleftmargin=0.5em,numbers=none,aboveskip=0.5em,belowskip=0.3em]
$( Unit Test 2: Missing whitespace between tokens $)
$( Should reject: True $)

$cwff$c|-.
$vx$.
wf$fwffx$.
ax$a|-x$.
\end{lstlisting}

\noindent\textbf{test06\_dangling\_dollar\_sign\_at\_eof.mm} [REJECT] \textit{L81-83: \$ must be followed by valid keyword char}\par

\begin{lstlisting}[language=Metamath,basicstyle=\ttfamily\scriptsize,frame=single,framerule=0.3pt,rulecolor=\color{black!20},backgroundcolor=\color{black!2},xleftmargin=1em,framexleftmargin=0.5em,numbers=none,aboveskip=0.5em,belowskip=0.3em]
$( Unit Test 6: Dangling dollar sign at EOF $)
$( Should reject: True $)

$c wff |- $.
$v x $.
wf $f wff x $.
$
\end{lstlisting}

\noindent\textbf{test18\_missing\_whitespace\_after\_comment.mm} [REJECT] \textit{L81-83: whitespace required between tokens}\par

\begin{lstlisting}[language=Metamath,basicstyle=\ttfamily\scriptsize,frame=single,framerule=0.3pt,rulecolor=\color{black!20},backgroundcolor=\color{black!2},xleftmargin=1em,framexleftmargin=0.5em,numbers=none,aboveskip=0.5em,belowskip=0.3em]
$( Unit Test 18: Missing whitespace after comment $)
$( Should reject: True $)

$c wff $.
$v x $.
$( No space after comment close $)wf $f wff x $.
\end{lstlisting}

\noindent\textbf{test25\_missing\_space\_before\_comment.mm} [REJECT] \textit{L81-83: \textbackslash{}\$( must be whitespace-separated}\par

\begin{lstlisting}[language=Metamath,basicstyle=\ttfamily\scriptsize,frame=single,framerule=0.3pt,rulecolor=\color{black!20},backgroundcolor=\color{black!2},xleftmargin=1em,framexleftmargin=0.5em,numbers=none,aboveskip=0.5em,belowskip=0.3em]
$( Unit Test 25: Missing space before comment $)
$( Should reject: True $)

$c wff$.$(comment$)
\end{lstlisting}

\noindent\textbf{test31\_missing\_whitespace\_after\_comment.mm} [REJECT] \textit{L81-83: \textbackslash{}\$) must be whitespace-separated}\par

\begin{lstlisting}[language=Metamath,basicstyle=\ttfamily\scriptsize,frame=single,framerule=0.3pt,rulecolor=\color{black!20},backgroundcolor=\color{black!2},xleftmargin=1em,framexleftmargin=0.5em,numbers=none,aboveskip=0.5em,belowskip=0.3em]
$( Unit Test 31: Missing whitespace after comment $)
$( Should reject: True $)

$( comment $)$c wff $.
\end{lstlisting}

\noindent\textbf{test49\_token\_splice\_axiom.mm} [REJECT] \textit{L81-83: token boundaries must be preserved}\par

\begin{lstlisting}[language=Metamath,basicstyle=\ttfamily\scriptsize,frame=single,framerule=0.3pt,rulecolor=\color{black!20},backgroundcolor=\color{black!2},xleftmargin=1em,framexleftmargin=0.5em,numbers=none,aboveskip=0.5em,belowskip=0.3em]
$( Unit Test 49: Include as token splice in axiom $)
$( Should reject: True - include inside statement is invalid in strict mode $)
$( Category: CORE - Strict include semantics $)
$( Spec Section 4.1.2: "include directives are processed at outermost level" $)
$( Spec Appendix E: Include not allowed inside statements $)

$c wff |- $.
$v x $.
fx $f wff x $.

$( Include splices tokens into the middle of this statement - INVALID $)
ax-splice $a $[ ./helpers/test49_helper.mm $]
\end{lstlisting}

\noindent\textbf{test50\_token\_splice\_proof.mm} [REJECT] \textit{L81-83: token boundaries must be preserved}\par

\begin{lstlisting}[language=Metamath,basicstyle=\ttfamily\scriptsize,frame=single,framerule=0.3pt,rulecolor=\color{black!20},backgroundcolor=\color{black!2},xleftmargin=1em,framexleftmargin=0.5em,numbers=none,aboveskip=0.5em,belowskip=0.3em]
$( Unit Test 50: Include as token splice in proof $)
$( Should reject: True - include inside statement is invalid in strict mode $)
$( Category: CORE - Strict include semantics $)
$( Spec Section 4.1.2: "include directives are processed at outermost level" $)
$( Spec Appendix E: Include not allowed inside statements $)

$c wff |- $.
$v x $.
fx $f wff x $.
ax-x $a |- x $.

$( Include splices proof steps from external file - INVALID $)
th $p |- x $= $[ ./helpers/test50_helper.mm $] $.
\end{lstlisting}

\subsection{Comment Delimiters (L89--91)}

\begin{quote}\small\itshape
4.1. SPECIFICATION OF THE METAMATH LANGUAGE
\end{quote}

\noindent\textbf{test03\_nested\_comment\_delimiters.mm} [REJECT] \textit{L89-91: \textbackslash{}\$( may not appear inside comments}\par

\begin{lstlisting}[language=Metamath,basicstyle=\ttfamily\scriptsize,frame=single,framerule=0.3pt,rulecolor=\color{black!20},backgroundcolor=\color{black!2},xleftmargin=1em,framexleftmargin=0.5em,numbers=none,aboveskip=0.5em,belowskip=0.3em]
$( Unit Test 3: Nested comment delimiters $)
$( Should reject: True $)

$c wff |- $.
$( outer $( nested $) comment $)
$v x $.
wf $f wff x $.
ax $a |- x $.
\end{lstlisting}

\subsection{Outermost Scope (L105--106)}

\begin{quote}\small\itshape
followed by \$]. It is only allowed in the outermost scope (i.e., not between \$\{
and \$\}) and must not be inside a statement (e.g., it may not occur between
\end{quote}

\noindent\textbf{test40\_include\_scope\_correct\_outer.mm} [REJECT] \textit{L105-106: \textbackslash{}\$[ \textbackslash{}\$] only allowed in outermost scope}\par

\begin{lstlisting}[language=Metamath,basicstyle=\ttfamily\scriptsize,frame=single,framerule=0.3pt,rulecolor=\color{black!20},backgroundcolor=\color{black!2},xleftmargin=1em,framexleftmargin=0.5em,numbers=none,aboveskip=0.5em,belowskip=0.3em]
$( Unit Test 40: Include inside block $)
$( Should reject: True - include must be at outermost level $)
$( Category: CORE - Strict include semantics $)
$( Spec Section 4.1.2: Include directives processed at outermost level only $)

$c wff |- $.
$v x $.
wx $f wff x $.
ax-x $a |- x $.

${
  $( Include inside block $)
  $[ ./helpers/test40_helper.mm $]

  $( Use included axiom INSIDE the block - this should work $)
  th1 $p |- y $= wy ax-inner $.
$}

$( Outside the block, only x is available $)
th2 $p |- x $= wx ax-x $.
\end{lstlisting}

\noindent\textbf{test47\_constant\_inner\_scope\_main.mm} [REJECT] \textit{L105-106: \textbackslash{}\$[ \textbackslash{}\$] only allowed in outermost scope}\par

\begin{lstlisting}[language=Metamath,basicstyle=\ttfamily\scriptsize,frame=single,framerule=0.3pt,rulecolor=\color{black!20},backgroundcolor=\color{black!2},xleftmargin=1em,framexleftmargin=0.5em,numbers=none,aboveskip=0.5em,belowskip=0.3em]
$( Unit Test 47: Include directive inside block $)
$( Category: CORE - CRITICAL scoping test $)
$( Should reject: True - include must be in outermost scope $)
$( Spec Section 4.1.2 L105-106: "$[ $] only allowed in outermost scope" $)

$c wff |- $.

${
  $( INVALID: Include inside block - must be in outermost scope only $)
  $[ ./helpers/test47_helper.mm $]

  $( Try to use included declarations $)
  th1 $p |- inner-var $= finner ax-inner $.
$}
\end{lstlisting}

\subsection{Block Delimiters (L153--155)}

\begin{quote}\small\itshape
typecode (an active constant), followed by an active variable, followed by
the \$. token. A \$e statement consists of a label, followed by \$e, followed
by its typecode (an active constant), followed by zero or more active math
\end{quote}

\noindent\textbf{test04\_unbalanced\_block\_delimiters.mm} [REJECT] \textit{L153-155: \textbackslash{}\$\{ must match \textbackslash{}\$\}}\par

\begin{lstlisting}[language=Metamath,basicstyle=\ttfamily\scriptsize,frame=single,framerule=0.3pt,rulecolor=\color{black!20},backgroundcolor=\color{black!2},xleftmargin=1em,framexleftmargin=0.5em,numbers=none,aboveskip=0.5em,belowskip=0.3em]
$( Unit Test 4: Unbalanced block delimiters $)
$( Should reject: True $)

$c wff |- $.
${
  $v x $.
  wf $f wff x $.
$( missing $} $)
ax $a |- x $.
\end{lstlisting}

\subsection{\$f Type Globality (L156--158)}

\begin{quote}\small\itshape
symbols, followed by the \$. token. A hypothesis is a \$f or \$e statement.
The type declared by a \$f statement for a given label is global even if the
variable is not (e.g., a database may not have wff P in one local scope nd
\end{quote}

\noindent\textbf{test15\_multiple\_f\_for\_same\_variable\_bad.mm} [REJECT] \textit{L156-158: type declared by \textbackslash{}\$f is GLOBAL (strict)}\par

\begin{lstlisting}[language=Metamath,basicstyle=\ttfamily\scriptsize,frame=single,framerule=0.3pt,rulecolor=\color{black!20},backgroundcolor=\color{black!2},xleftmargin=1em,framexleftmargin=0.5em,numbers=none,aboveskip=0.5em,belowskip=0.3em]
$( Unit Test 15: Multiple $f for same variable $)
$( Should reject: True $)

$c wff class $.
$v x $.
f1 $f wff x $.
f2 $f class x $.
\end{lstlisting}

\noindent\textbf{test16\_conflicting\_typecodes\_bad.mm} [REJECT] \textit{L156-158: type declared by \textbackslash{}\$f is GLOBAL (strict)}\par

\begin{lstlisting}[language=Metamath,basicstyle=\ttfamily\scriptsize,frame=single,framerule=0.3pt,rulecolor=\color{black!20},backgroundcolor=\color{black!2},xleftmargin=1em,framexleftmargin=0.5em,numbers=none,aboveskip=0.5em,belowskip=0.3em]
$( Unit Test 16: Conflicting typecodes $)
$( Should reject: True $)

$c wff class $.
$v x $.
${
  f1 $f wff x $.
  ${
    f2 $f class x $.
  $}
$}
\end{lstlisting}

\subsection{Constant/Variable Declarations (L161--163)}

\begin{quote}\small\itshape
variables, followed by the \$. token. A compound \$d statement consists
of \$d, followed by three or more variables (all different), followed by the
\$. token. The order of the variables in a \$d statement is unimportant. A
\end{quote}

\noindent\textbf{local-var-bad.mm} [REJECT] \textit{L161-163: variable scope violation}\par

\begin{lstlisting}[language=Metamath,basicstyle=\ttfamily\scriptsize,frame=single,framerule=0.3pt,rulecolor=\color{black!20},backgroundcolor=\color{black!2},xleftmargin=1em,framexleftmargin=0.5em,numbers=none,aboveskip=0.5em,belowskip=0.3em]
  $c formula $.
  $c A. $.
  $c setvar $.
  $v ph $.
  wph $f formula ph $.

  ${
    $v x $.
    $( This is not valid because there is no $f for the variable ` x ` $)
    $( Extend formula definition to include the universal quantifier ('for all').  $)
    wal $a formula A. x ph $.
  $}
\end{lstlisting}

\noindent\textbf{test07\_redeclaration\_of\_constant.mm} [REJECT] \textit{L161-163: constant may not be redeclared}\par

\begin{lstlisting}[language=Metamath,basicstyle=\ttfamily\scriptsize,frame=single,framerule=0.3pt,rulecolor=\color{black!20},backgroundcolor=\color{black!2},xleftmargin=1em,framexleftmargin=0.5em,numbers=none,aboveskip=0.5em,belowskip=0.3em]
$( Unit Test 7: Redeclaration of constant $)
$( Should reject: True $)

$c wff |- $.
$c wff $.
$v x $.
wf $f wff x $.
ax $a |- x $.
\end{lstlisting}

\noindent\textbf{test45\_variable\_redeclaration.mm} [ACCEPT] \textit{L161-163: variables MAY be redeclared in new scope}\par

\begin{lstlisting}[language=Metamath,basicstyle=\ttfamily\scriptsize,frame=single,framerule=0.3pt,rulecolor=\color{black!20},backgroundcolor=\color{black!2},xleftmargin=1em,framexleftmargin=0.5em,numbers=none,aboveskip=0.5em,belowskip=0.3em]
$( Unit Test 45: Variable redeclaration across scopes $)
$( Should reject: False - variable can be redeclared after scope ends $)
$( Spec Section 4.2.8: "A variable may be declared again after it becomes inactive" $)

$c wff class |- $.
$v x $.

${
  $( First scope: x has typecode wff $)
  fx1 $f wff x $.
  ax1 $a |- x $.
  $( Use x within its scope $)
  th1 $p |- x $= fx1 ax1 $.
$}

${
  $( Second scope: x has typecode class - VALID redeclaration $)
  fx2 $f class x $.
  ax2 $a class x $.
  $( Use x within its scope $)
  th2 $p class x $= fx2 ax2 $.
$}
\end{lstlisting}

\noindent\textbf{test48\_variable\_conflict\_main.mm} [REJECT] \textit{L161-163: variable conflicts with constant}\par

\begin{lstlisting}[language=Metamath,basicstyle=\ttfamily\scriptsize,frame=single,framerule=0.3pt,rulecolor=\color{black!20},backgroundcolor=\color{black!2},xleftmargin=1em,framexleftmargin=0.5em,numbers=none,aboveskip=0.5em,belowskip=0.3em]
$( Unit Test 48: Include causing variable redeclaration conflict $)
$( Category: CORE - CRITICAL scoping test $)
$( Should reject: True - x already active when include tries to redeclare it $)
$( Spec Section 4.2.8: "A variable may not be declared a second time while it is active" $)

$c wff |- $.
$v x $.
fx $f wff x $.

$( Include file that also declares $v x - should fail $)
$[ ./helpers/test48_helper.mm $]
\end{lstlisting}

\subsection{Label Syntax (L167--169)}

\begin{quote}\small\itshape
A \$d statement is also called a disjoint (or distinct) variable restriction.
A \$a statement consists of a label, followed by \$a, followed by its
typecode (an active constant), followed by zero or more active math symbols,
\end{quote}

\noindent\textbf{demo0-illegal-label-bad1.mm} [REJECT] \textit{L167-169: labels must be alphanumeric + hyphen + underscore + period}\par

\begin{lstlisting}[language=Metamath,basicstyle=\ttfamily\scriptsize,frame=single,framerule=0.3pt,rulecolor=\color{black!20},backgroundcolor=\color{black!2},xleftmargin=1em,framexleftmargin=0.5em,numbers=none,aboveskip=0.5em,belowskip=0.3em]
$( demo0.mm  1-Jan-04 $)

$(
                      PUBLIC DOMAIN DEDICATION

This file is placed in the public domain per the Creative Commons Public
Domain Dedication. ht\tp://creativecommons.org/licenses/publicdomain/

Norman Megill - email: nm at alum.mit.edu

$)


$( This file is the introductory formal system example described
   in Chapter 2 of the Metamath book. $)

$( Declare the constant symbols we will use $)
    $c 0 + = -> ( ) term wff |- $.
$( Declare the metavariables we will use $)
    $v t r s P Q $.
$( Specify properties of the metavariables $)
    t\t $f term t $.
    tr $f term r $.
    ts $f term s $.
    wp $f wff P $.
    wq $f wff Q $.
$( Define "term" (part 1) $)
    (tze $a term 0 $.
$( Define "term" (part 2) $)
    t)pl $a term ( t + r ) $.
$( Define "wff" (part 1) $)
    w"eq $a wff t = r $.
$( Define "wff" (part 2) $)
    w%i\m $a wff ( P -> Q ) $.
$( State axiom a1 $)
    a1 $a |- ( t = r -> ( t = s -> r = s ) ) $.
$( State axiom 'a2` $)
    'a2` $a |- ( t + 0 ) = t $.
    ${
       min $e |- P $.
       maj $e |- ( P -> Q ) $.
$( Define the modus ponens inference rule $)
       mp  $a |- Q $.
    $}
$( Prove a theorem $)
    th1 $p |- t = t $=
  $( Here is its proof: $)
       t\t (tze t)pl t\t w"eq t\t t\t w"eq t\t 'a2` t\t (tze t)pl
       t\t w"eq t\t (tze t)pl t\t w"eq t\t t\t w"eq w%i\m t\t 'a2`
       t\t (tze t)pl t\t t\t a1 mp mp
     $.
\end{lstlisting}

\noindent\textbf{underscores.mm} [ACCEPT] \textit{L167-169: underscores allowed in labels}\par

\begin{lstlisting}[language=Metamath,basicstyle=\ttfamily\scriptsize,frame=single,framerule=0.3pt,rulecolor=\color{black!20},backgroundcolor=\color{black!2},xleftmargin=1em,framexleftmargin=0.5em,numbers=none,aboveskip=0.5em,belowskip=0.3em]
  $c a $.

  $( This is _italic_ and the command "MM-PA> MINIMIZE__WITH *" and a
     sub_script, a label ~ lab[[``~~__ , and a url: ~ http://a_b__c.com/[[``~~__.
     Doubled chars: [[ `` ~~ __
     <HTML>
     Inside HTML tags:

     This is _italic_ and the command "MM-PA> MINIMIZE__WITH *" and a
     sub_script, a label ~ lab[[``~~__ , and a url: ~ http://a_b__c.com[[``~~__.
     Doubled chars: [[ `` ~~ __
     </HTML> $)
  underscores $a a $.

$( $t
 htmldef "a" as "a";
  althtmldef "a" as "a";
   latexdef "a" as "a"; $)
\end{lstlisting}

\noindent\textbf{test57\_case\_sensitivity\_labels.mm} [ACCEPT] \textit{L167-169: labels are case-sensitive}\par

\begin{lstlisting}[language=Metamath,basicstyle=\ttfamily\scriptsize,frame=single,framerule=0.3pt,rulecolor=\color{black!20},backgroundcolor=\color{black!2},xleftmargin=1em,framexleftmargin=0.5em,numbers=none,aboveskip=0.5em,belowskip=0.3em]
$( Test 57: Case sensitivity for labels - labels are case-sensitive per Sec 4.1.3; should ACCEPT $)

$c wff |- $.
$v ph $.

f1 $f wff ph $.
ax $a |- ph $.
Ax $a |- ph $.

p1 $p |- ph $= f1 ax $.
p2 $p |- ph $= f1 Ax $.
\end{lstlisting}

\subsection{\$f Scope (L174--175)}

\begin{quote}\small\itshape
A \$f, \$e, or \$d statement is active from the place it occurs until the
end of the block it occurs in. A \$a or \$p statement is active from the place
\end{quote}

\noindent\textbf{test35\_f\_not\_active\_after\_block\_close.mm} [REJECT] \textit{L174-175: \textbackslash{}\$f active only until block end}\par

\begin{lstlisting}[language=Metamath,basicstyle=\ttfamily\scriptsize,frame=single,framerule=0.3pt,rulecolor=\color{black!20},backgroundcolor=\color{black!2},xleftmargin=1em,framexleftmargin=0.5em,numbers=none,aboveskip=0.5em,belowskip=0.3em]
$( Unit Test 35: $f not active after block close $)
$( Should reject: True $)

$c wff |- $.
${
  $v x $.
  wf $f wff x $.
$}
$( x and wf are no longer active $)
bad $a |- x $.
\end{lstlisting}

\subsection{Variable Scope (L174--178)}

\begin{quote}\small\itshape
A \$f, \$e, or \$d statement is active from the place it occurs until the
end of the block it occurs in. A \$a or \$p statement is active from the place
it occurs through the end of the database. There may not be two active \$f
statements containing the same variable. Each variable in a \$e, \$a, or \$p
statement must exist in an active \$f statement.3
\end{quote}

\noindent\textbf{test08\_variable\_used\_outside\_scope.mm} [REJECT] \textit{L174-178: variable must have active \textbackslash{}\$f}\par

\begin{lstlisting}[language=Metamath,basicstyle=\ttfamily\scriptsize,frame=single,framerule=0.3pt,rulecolor=\color{black!20},backgroundcolor=\color{black!2},xleftmargin=1em,framexleftmargin=0.5em,numbers=none,aboveskip=0.5em,belowskip=0.3em]
$( Unit Test 8: Variable used outside scope $)
$( Should reject: True $)

$c wff |- $.
${
  $v x $.
  wf $f wff x $.
$}
bad $a |- x $.
\end{lstlisting}

\noindent\textbf{test13\_f\_with\_undeclared\_variable.mm} [REJECT] \textit{L174-178: \textbackslash{}\$f variable must be declared}\par

\begin{lstlisting}[language=Metamath,basicstyle=\ttfamily\scriptsize,frame=single,framerule=0.3pt,rulecolor=\color{black!20},backgroundcolor=\color{black!2},xleftmargin=1em,framexleftmargin=0.5em,numbers=none,aboveskip=0.5em,belowskip=0.3em]
$( Unit Test 13: $f with undeclared variable $)
$( Should reject: True $)

$c wff |- $.
bad $f wff undeclared $.
\end{lstlisting}

\subsection{Variable Must Have \$f (L177--178)}

\begin{quote}\small\itshape
statements containing the same variable. Each variable in a \$e, \$a, or \$p
statement must exist in an active \$f statement.3
\end{quote}

\noindent\textbf{test14\_variable\_without\_f\_hypothesis.mm} [REJECT] \textit{L177-178: each variable in \textbackslash{}\$e/\textbackslash{}\$a/\textbackslash{}\$p must have active \textbackslash{}\$f}\par

\begin{lstlisting}[language=Metamath,basicstyle=\ttfamily\scriptsize,frame=single,framerule=0.3pt,rulecolor=\color{black!20},backgroundcolor=\color{black!2},xleftmargin=1em,framexleftmargin=0.5em,numbers=none,aboveskip=0.5em,belowskip=0.3em]
$( Unit Test 14: Variable without $f hypothesis $)
$( Should reject: True $)

$c wff |- $.
$v x $.
bad $a |- x $.
\end{lstlisting}

\subsection{Label Uniqueness (L179--180)}

\begin{quote}\small\itshape
Each label token must be unique, and no label token may match any
math symbol token.4
\end{quote}

\noindent\textbf{test10\_duplicate\_labels.mm} [REJECT] \textit{L179-180: label tokens must be unique}\par

\begin{lstlisting}[language=Metamath,basicstyle=\ttfamily\scriptsize,frame=single,framerule=0.3pt,rulecolor=\color{black!20},backgroundcolor=\color{black!2},xleftmargin=1em,framexleftmargin=0.5em,numbers=none,aboveskip=0.5em,belowskip=0.3em]
$( Unit Test 10: Duplicate labels $)
$( Should reject: True $)

$c wff |- $.
$v x $.
dup $f wff x $.
dup $a |- x $.
\end{lstlisting}

\noindent\textbf{test11\_label\_conflicts\_with\_symbol.mm} [REJECT] \textit{L179-180: no label may match any math symbol}\par

\begin{lstlisting}[language=Metamath,basicstyle=\ttfamily\scriptsize,frame=single,framerule=0.3pt,rulecolor=\color{black!20},backgroundcolor=\color{black!2},xleftmargin=1em,framexleftmargin=0.5em,numbers=none,aboveskip=0.5em,belowskip=0.3em]
$( Unit Test 11: Label conflicts with symbol $)
$( Should reject: True $)

$c wff |- $.
$v x $.
wff $f wff x $.
\end{lstlisting}

\noindent\textbf{test34\_label\_token\_in\_math\_context.mm} [REJECT] \textit{L179-180,339: labels and math symbols are disjoint}\par

\begin{lstlisting}[language=Metamath,basicstyle=\ttfamily\scriptsize,frame=single,framerule=0.3pt,rulecolor=\color{black!20},backgroundcolor=\color{black!2},xleftmargin=1em,framexleftmargin=0.5em,numbers=none,aboveskip=0.5em,belowskip=0.3em]
$( Unit Test 34: Label token in math context $)
$( Should reject: True $)

$c wff |- $.
$v x $.
wf $f wff x $.
$( Using 'wf' where math symbol expected $)
bad $a wf x $.
\end{lstlisting}

\subsection{Mandatory Hypotheses (L182--184)}

\begin{quote}\small\itshape
of (zero or more) variables in the assertion and in any active \$e statements.
The (possibly empty) set of mandatory hypotheses is the set of all active
\$f statements containing mandatory variables, together with all active \$e
\end{quote}

\noindent\textbf{test64\_multiple\_ehyp\_order.mm} [ACCEPT] \textit{L182-184, L211-214: mandatory hyp ordering}\par

\begin{lstlisting}[language=Metamath,basicstyle=\ttfamily\scriptsize,frame=single,framerule=0.3pt,rulecolor=\color{black!20},backgroundcolor=\color{black!2},xleftmargin=1em,framexleftmargin=0.5em,numbers=none,aboveskip=0.5em,belowskip=0.3em]
$( test64_multiple_ehyp_order.mm - ACCEPT
   Tests spec 4.1.4 L182-184, L211-214:
   "The (possibly empty) set of mandatory hypotheses is the set of all
   active $f statements containing mandatory variables, together with
   all active $e statements."

   "causes them to match the topmost entries of the stack, in order of
   occurrence of the mandatory hypotheses"

   Key insight: mandatory hypotheses = $f (for mandatory vars) + $e (all active)
   in order of OCCURRENCE in the file.
$)

$c wff |- ( ) -> $.
$v p q $.
wp $f wff p $.
wq $f wff q $.

$( Build implication $)
wi $a wff ( p -> q ) $.

$( A theorem with two essential hypotheses.
   Mandatory hyps in occurrence order: wp, wq, e1, e2 $)
${
  e1 $e |- p $.
  e2 $e |- ( p -> q ) $.
  $( Modus ponens: from p and p->q, derive q $)
  mp $a |- q $.
$}

$( Use mp with concrete values.
   To apply mp where p:=( p -> p ), q:=( p -> p ):
   Push order: wp (for p), wp (for q), inst of e1, inst of e2

   We'll prove: from ( p -> p ) and ( ( p -> p ) -> ( p -> p ) ), get ( p -> p )
$)

${
  h1 $e |- ( p -> p ) $.
  h2 $e |- ( ( p -> p ) -> ( p -> p ) ) $.

  $( Stack for mp: wp wp wi, wp wp wi, h1, h2
     This substitutes p:=( p -> p ), q:=( p -> p ) $)
  result $p |- ( p -> p ) $=
    wp wp wi wp wp wi h1 h2 mp $.
$}
\end{lstlisting}

\subsection{Self-Substitution (L204--205)}

\begin{quote}\small\itshape
An expression is a sequence of math symbols. A substitution map
associates a set of variables with a set of expressions. It is acceptable for a
\end{quote}

\noindent\textbf{test62\_self\_substitution.mm} [ACCEPT] \textit{L204-205: variable mapped to expression containing itself}\par

\begin{lstlisting}[language=Metamath,basicstyle=\ttfamily\scriptsize,frame=single,framerule=0.3pt,rulecolor=\color{black!20},backgroundcolor=\color{black!2},xleftmargin=1em,framexleftmargin=0.5em,numbers=none,aboveskip=0.5em,belowskip=0.3em]
$( test62_self_substitution.mm - ACCEPT
   Tests spec 4.1.4 L204-205:
   "It is acceptable for a variable to be mapped to an expression
   containing it."

   This tests that p can be substituted with an expression containing p.
$)

$c wff |- ( ) -> $.
$v p q $.
wp $f wff p $.
wq $f wff q $.

$( Build implication: wff ( p -> q ) $)
wi $a wff ( p -> q ) $.

$( Axiom: |- ( p -> p )
   Mandatory hypothesis: just wp (only var p is used) $)
ax-id $a |- ( p -> p ) $.

$( Prove ( ( p -> p ) -> ( p -> p ) )
   We apply ax-id with p := ( p -> p )
   This means p maps to an expression containing p itself!

   ax-id needs: wp
   We need to push wff ( p -> p ) to substitute for p.
   Building wff ( p -> p ): wp wp wi (gives wff ( p -> p ))
   Then apply ax-id.
$)
th-self $p |- ( ( p -> p ) -> ( p -> p ) ) $=
  wp wp wi ax-id $.
\end{lstlisting}

\subsection{Typecode Matching (L208--212)}

\begin{quote}\small\itshape
the expressions that the variables map to.
A proof is scanned in order of its label sequence. If the label refers to an
active hypothesis, the expression in the hypothesis is pushed onto a stack.
If the label refers to an assertion, a (unique) substitution must exist that,
when made to the mandatory hypotheses of the referenced assertion, causes
\end{quote}

\noindent\textbf{test22\_typecode\_mismatch\_in\_substitution.mm} [REJECT] \textit{L208-212: substitution must match typecode}\par

\begin{lstlisting}[language=Metamath,basicstyle=\ttfamily\scriptsize,frame=single,framerule=0.3pt,rulecolor=\color{black!20},backgroundcolor=\color{black!2},xleftmargin=1em,framexleftmargin=0.5em,numbers=none,aboveskip=0.5em,belowskip=0.3em]
$( Unit Test 22: Typecode mismatch in substitution $)
$( Should reject: True $)

$c wff class |- $.
$v x y $.
f1 $f wff x $.
f2 $f class y $.
ax $a |- x $.
bad $p |- y $= f2 ax $.
\end{lstlisting}

\subsection{Proof Labels (L213--218)}

\begin{quote}\small\itshape
them to match the topmost (i.e. most recent) entries of the stack, in order
of occurrence of the mandatory hypotheses, with the topmost stack entry
matching the last mandatory hypothesis of the referenced assertion. As many
stack entries as there are mandatory hypotheses are then popped from the
stack. The same substitution is made to the referenced assertion, and the
result is pushed onto the stack. After the last label in the proof is processed,
\end{quote}

\noindent\textbf{test21\_self-referential\_proof.mm} [REJECT] \textit{L213-218: proof may not reference itself}\par

\begin{lstlisting}[language=Metamath,basicstyle=\ttfamily\scriptsize,frame=single,framerule=0.3pt,rulecolor=\color{black!20},backgroundcolor=\color{black!2},xleftmargin=1em,framexleftmargin=0.5em,numbers=none,aboveskip=0.5em,belowskip=0.3em]
$( Unit Test 21: Self-referential proof $)
$( Should reject: True $)

$c wff |- -> ( ) $.
$v x $.
wf $f wff x $.
evil $p |- ( x -> x ) $= ( evil ) A $.
\end{lstlisting}

\noindent\textbf{test24\_undefined\_label\_in\_proof.mm} [REJECT] \textit{L213-218: proof labels must be defined}\par

\begin{lstlisting}[language=Metamath,basicstyle=\ttfamily\scriptsize,frame=single,framerule=0.3pt,rulecolor=\color{black!20},backgroundcolor=\color{black!2},xleftmargin=1em,framexleftmargin=0.5em,numbers=none,aboveskip=0.5em,belowskip=0.3em]
$( Unit Test 24: Undefined label in proof $)
$( Should reject: True $)

$c wff |- $.
$v x $.
wf $f wff x $.
ax $a |- x $.
bad $p |- x $= nosuchlabel $.
\end{lstlisting}

\noindent\textbf{test32\_forward\_reference\_in\_proof.mm} [REJECT] \textit{L213-218: proof labels must be previously defined}\par

\begin{lstlisting}[language=Metamath,basicstyle=\ttfamily\scriptsize,frame=single,framerule=0.3pt,rulecolor=\color{black!20},backgroundcolor=\color{black!2},xleftmargin=1em,framexleftmargin=0.5em,numbers=none,aboveskip=0.5em,belowskip=0.3em]
$( Unit Test 32: Forward reference in proof $)
$( Should reject: True $)

$c wff |- $.
$v x $.
wf $f wff x $.
early $p |- x $= wf later $.
later $a |- x $.
\end{lstlisting}

\subsection{Stack Matching (L218--220)}

\begin{quote}\small\itshape
result is pushed onto the stack. After the last label in the proof is processed,
the stack must have a single entry that matches the expression in the \$p
statement containing the proof.
\end{quote}

\noindent\textbf{test26\_wrong\_conclusion\_in\_proof.mm} [REJECT] \textit{L218-220: final stack must match \textbackslash{}\$p statement}\par

\begin{lstlisting}[language=Metamath,basicstyle=\ttfamily\scriptsize,frame=single,framerule=0.3pt,rulecolor=\color{black!20},backgroundcolor=\color{black!2},xleftmargin=1em,framexleftmargin=0.5em,numbers=none,aboveskip=0.5em,belowskip=0.3em]
$( Unit Test 26: Wrong conclusion in proof $)
$( Should reject: True $)

$c wff |- $.
$v x y $.
f1 $f wff x $.
f2 $f wff y $.
ax $a |- x $.
bad $p |- y $= f1 ax $.
\end{lstlisting}

\noindent\textbf{test33\_compressed\_proof\_stack\_underflow.mm} [REJECT] \textit{L218-220: stack must not underflow}\par

\begin{lstlisting}[language=Metamath,basicstyle=\ttfamily\scriptsize,frame=single,framerule=0.3pt,rulecolor=\color{black!20},backgroundcolor=\color{black!2},xleftmargin=1em,framexleftmargin=0.5em,numbers=none,aboveskip=0.5em,belowskip=0.3em]
$( Unit Test 33: Compressed proof stack underflow $)
$( Should reject: True $)

$c wff |- $.
$v x $.
wf $f wff x $.
ax $a |- x $.
$( C tries to pop 3rd item, but stack has only 1 $)
bad $p |- x $= ( ax ) C $.
\end{lstlisting}

\subsection{Unknown Steps (L221--223)}

\begin{quote}\small\itshape
A proof may contain a ? in place of a label to indicate an unknown step
(Section 4.4.6). A proof verifier may ignore any proof containing ? but should
warn the user that the proof is incomplete.
\end{quote}

\subsection{Typecodes in \$f (L286--292)}

\begin{quote}\small\itshape
Constant symbol declaration
Variable symbol declaration
Disjoint variable restriction
Variable-type (“floating”) hypothesis
Logical (“essential”) hypothesis
Axiomatic assertion
Provable assertion
\end{quote}

\noindent\textbf{test12\_non-constant\_typecode.mm} [REJECT] \textit{L286-292: first symbol in \textbackslash{}\$f must be constant (typecode)}\par

\begin{lstlisting}[language=Metamath,basicstyle=\ttfamily\scriptsize,frame=single,framerule=0.3pt,rulecolor=\color{black!20},backgroundcolor=\color{black!2},xleftmargin=1em,framexleftmargin=0.5em,numbers=none,aboveskip=0.5em,belowskip=0.3em]
$( Unit Test 12: Non-constant typecode $)
$( Should reject: True $)

$c wff |- $.
$v x y $.
bad $f x y $.
\end{lstlisting}

\subsection{Math Symbol Syntax (L339--342)}

\begin{quote}\small\itshape
Math symbols are tokens used to represent the symbols that appear in
ordinary mathematical formulas. They may consist of any combination
of the 93 non-whitespace printable ascii characters other than \$ . Some
examples are x, +, (, |-, !\%@?\&, and bounded. For readability, it is best
\end{quote}

\noindent\textbf{demo0-dollar-in-math-symbol-bad1.mm} [REJECT] \textit{L339-342: math symbols must not contain \$}\par

\begin{lstlisting}[language=Metamath,basicstyle=\ttfamily\scriptsize,frame=single,framerule=0.3pt,rulecolor=\color{black!20},backgroundcolor=\color{black!2},xleftmargin=1em,framexleftmargin=0.5em,numbers=none,aboveskip=0.5em,belowskip=0.3em]
$( demo0.mm  1-Jan-04 $)

$(
                      PUBLIC DOMAIN DEDICATION

This file is placed in the public domain per the Creative Commons Public
Domain Dedication. http://creativecommons.org/licenses/publicdomain/

Norman Megill - email: nm at alum.mit.edu

$)


$( This file is the introductory formal system example described
   in Chapter 2 of the Meamath book. $)

$( Declare the constant symbols we will use $)
    $c 0 + = -> ( ) meep$ term wff |- $.
$( Declare the metavariables we will use $)
    $v t r s P Q $.
\end{lstlisting}

\noindent\textbf{test05\_dollar\_sign\_in\_math\_symbols.mm} [REJECT] \textit{L339-342: math symbols must not contain \$}\par

\begin{lstlisting}[language=Metamath,basicstyle=\ttfamily\scriptsize,frame=single,framerule=0.3pt,rulecolor=\color{black!20},backgroundcolor=\color{black!2},xleftmargin=1em,framexleftmargin=0.5em,numbers=none,aboveskip=0.5em,belowskip=0.3em]
$( Unit Test 5: Dollar sign in math symbols $)
$( Should reject: True $)

$c wff |- $.
$c a$b $.
$v x $.
wf $f wff x $.
ax $a |- x $.
\end{lstlisting}

\noindent\textbf{test58\_case\_sensitivity\_symbols.mm} [ACCEPT] \textit{L339-342: math symbols are case-sensitive}\par

\begin{lstlisting}[language=Metamath,basicstyle=\ttfamily\scriptsize,frame=single,framerule=0.3pt,rulecolor=\color{black!20},backgroundcolor=\color{black!2},xleftmargin=1em,framexleftmargin=0.5em,numbers=none,aboveskip=0.5em,belowskip=0.3em]
$( Test 58: Case sensitivity for symbols - constants are case-sensitive per Sec 4.1.3; should ACCEPT $)

$c WFF PH |- $.
$v ph $.

f1 $f WFF ph $.
a1 $a |- ph $.
p1 $p |- ph $= f1 a1 $.
\end{lstlisting}

\subsection{\$d Syntax (L547--549)}

\begin{quote}\small\itshape
for each possible pair of variables in the original \$d statement. For example,
\$d w x y z \$.
is equivalent to
\end{quote}

\noindent\textbf{test09\_d\_with\_non-variables.mm} [REJECT] \textit{L547-549: \textbackslash{}\$d must contain only variables}\par

\begin{lstlisting}[language=Metamath,basicstyle=\ttfamily\scriptsize,frame=single,framerule=0.3pt,rulecolor=\color{black!20},backgroundcolor=\color{black!2},xleftmargin=1em,framexleftmargin=0.5em,numbers=none,aboveskip=0.5em,belowskip=0.3em]
$( Unit Test 9: $d with non-variables $)
$( Should reject: True $)

$c wff |- $.
$d wff |- $.
\end{lstlisting}

\subsection{Duplicate \$d (L553--558)}

\begin{quote}\small\itshape
\$d x y \$.
\$d x z \$.
\$d y z \$.
Two or more simple \$d statements specifying the same variable pair
are internally combined into a single \$d statement. Thus the set of three
statements
\end{quote}

\noindent\textbf{test56\_duplicate\_d\_normalization\_bad.mm} [REJECT] \textit{L553-558: duplicate \textbackslash{}\$d is error}\par

\begin{lstlisting}[language=Metamath,basicstyle=\ttfamily\scriptsize,frame=single,framerule=0.3pt,rulecolor=\color{black!20},backgroundcolor=\color{black!2},xleftmargin=1em,framexleftmargin=0.5em,numbers=none,aboveskip=0.5em,belowskip=0.3em]
$( Test 56: Duplicate $d normalization - duplicates allowed; should ACCEPT per Sec 4.2.4 (union of DV pairs) $)

$c wff |- $.
$v x $.

$d x x $.
$d x x $.

f1 $f wff x $.
a1 $a |- x $.
p1 $p |- x $= f1 a1 $.
\end{lstlisting}

\subsection{\$c Outermost Scope (L1088)}

\begin{quote}\small\itshape
All \$c statements must be placed in the outermost block. Since they are
\end{quote}

\noindent\textbf{test47b\_constant\_inner\_scope\_direct.mm} [REJECT] \textit{L1088: \textbackslash{}\$c must be in outermost block}\par

\begin{lstlisting}[language=Metamath,basicstyle=\ttfamily\scriptsize,frame=single,framerule=0.3pt,rulecolor=\color{black!20},backgroundcolor=\color{black!2},xleftmargin=1em,framexleftmargin=0.5em,numbers=none,aboveskip=0.5em,belowskip=0.3em]
$( Unit Test 47b: $c declaration directly in inner scope $)
$( Category: CORE - CRITICAL scoping test $)
$( Should reject: True - $c must be in outermost block $)
$( Spec Section 4.1.2 L1088: "All $c statements must be placed in the outermost block" $)

$c wff |- $.

${
  $( INVALID: $c declaration inside block - must be in outermost scope only $)
  $c inner-const $.

  $v x $.
  wx $f wff x $.
  ax-x $a |- x $.
$}
\end{lstlisting}

\subsection{Missing \$p (L1114--1115)}

\begin{quote}\small\itshape
A constant may not be redeclared. And, as mentioned above, constants
must be declared in the outermost block.
\end{quote}

\noindent\textbf{missing-dollar-p.mm} [REJECT] \textit{L1114-1115: missing \textbackslash{}\$p statement}\par

\begin{lstlisting}[language=Metamath,basicstyle=\ttfamily\scriptsize,frame=single,framerule=0.3pt,rulecolor=\color{black!20},backgroundcolor=\color{black!2},xleftmargin=1em,framexleftmargin=0.5em,numbers=none,aboveskip=0.5em,belowskip=0.3em]
${ mytoken $}
\end{lstlisting}

\subsection{Compressed Proof ? (L1894--1895)}

\begin{quote}\small\itshape
if desired. In this case, you will see one or more ?’s in the compressed-proof
part.
\end{quote}

\noindent\textbf{test30\_qmark\_in\_compressed\_proof.mm} [ACCEPT] \textit{L1894-1895: ? allowed in compressed proofs}\par

\begin{lstlisting}[language=Metamath,basicstyle=\ttfamily\scriptsize,frame=single,framerule=0.3pt,rulecolor=\color{black!20},backgroundcolor=\color{black!2},xleftmargin=1em,framexleftmargin=0.5em,numbers=none,aboveskip=0.5em,belowskip=0.3em]
$( Unit Test 30: ? in compressed proof $)
$( Should reject: False $)

$c wff |- $.
$v x $.
wf $f wff x $.
ax $a |- x $.
incomplete $p |- x $= ( ax ) ? $.
\end{lstlisting}

\subsection{Proof Verification Failures (\S4.2.1)}

\noindent\textbf{anatomy-bad1.mm} [REJECT] \textit{4.2.1: proof verification failure}\par

\begin{lstlisting}[language=Metamath,basicstyle=\ttfamily\scriptsize,frame=single,framerule=0.3pt,rulecolor=\color{black!20},backgroundcolor=\color{black!2},xleftmargin=1em,framexleftmargin=0.5em,numbers=none,aboveskip=0.5em,belowskip=0.3em]
$( anatomy.mm $)

$(
                      PUBLIC DOMAIN DEDICATION

This file is placed in the public domain per the Creative Commons Public
Domain Dedication. http://creativecommons.org/licenses/publicdomain/

Norman Megill - email: nm at alum.mit.edu

$)


$( This file is the introductory formal system example described
   in Chapter 4 of the Metamath book, "Anatomy of a proof". $)

$( Declare the constant symbols we will use $)
$c ( ) -> wff $.
$( Declare the metavariables we will use $)
$v p q r s $.
$( Declare wffs $)
wp $f wff p $.
wq $f wff q $.
wr $f wff r $.
ws $f wff s $.
w2 $a wff ( p -> q ) $.

$( Prove a simple theorem. $)
wnew $p wff ( s -> ( r -> p ) ) $= ws wr wp w2 ws $.
\end{lstlisting}

\noindent\textbf{anatomy-bad2.mm} [REJECT] \textit{4.2.1: proof verification failure}\par

\begin{lstlisting}[language=Metamath,basicstyle=\ttfamily\scriptsize,frame=single,framerule=0.3pt,rulecolor=\color{black!20},backgroundcolor=\color{black!2},xleftmargin=1em,framexleftmargin=0.5em,numbers=none,aboveskip=0.5em,belowskip=0.3em]
$( anatomy.mm $)

$(
                      PUBLIC DOMAIN DEDICATION

This file is placed in the public domain per the Creative Commons Public
Domain Dedication. http://creativecommons.org/licenses/publicdomain/

Norman Megill - email: nm at alum.mit.edu

$)


$( This file is the introductory formal system example described
   in Chapter 4 of the Metamath book, "Anatomy of a proof". $)

$( Declare the constant symbols we will use $)
$c ( ) -> wff $.
$( Declare the metavariables we will use $)
$v p q r s $.
$( Declare wffs $)
wp $f wff p $.
wq $f wff q $.
wr $f wff r $.
ws $f wff s $.
w2 $a wff ( p -> q ) $.

$( Prove a simple theorem. $)
wnew $p wff ( s -> ( r -> p ) ) $= ws wr wp w2 $.
\end{lstlisting}

\noindent\textbf{anatomy-bad3.mm} [REJECT] \textit{4.2.1: proof verification failure}\par

\begin{lstlisting}[language=Metamath,basicstyle=\ttfamily\scriptsize,frame=single,framerule=0.3pt,rulecolor=\color{black!20},backgroundcolor=\color{black!2},xleftmargin=1em,framexleftmargin=0.5em,numbers=none,aboveskip=0.5em,belowskip=0.3em]
$( anatomy.mm $)

$(
                      PUBLIC DOMAIN DEDICATION

This file is placed in the public domain per the Creative Commons Public
Domain Dedication. http://creativecommons.org/licenses/publicdomain/

Norman Megill - email: nm at alum.mit.edu

$)


$( This file is the introductory formal system example described
   in Chapter 4 of the Metamath book, "Anatomy of a proof". $)

$( Declare the constant symbols we will use $)
$c ( ) -> wff $.
$( Declare the metavariables we will use $)
$v p q r s $.
$( Declare wffs $)
wp $f wff p $.
wq $f wff q $.
wr $f wff r $.
ws $f wff s $.
w2 $a wff ( p -> q ) $.

$( Prove a simple theorem. $)
wnew $p wff ( s -> ( r -> p ) ) $= wr wp w2 w2 $.
\end{lstlisting}

\noindent\textbf{demo0-bad1.mm} [REJECT] \textit{4.2.1: proof verification failure}\par

\begin{lstlisting}[language=Metamath,basicstyle=\ttfamily\scriptsize,frame=single,framerule=0.3pt,rulecolor=\color{black!20},backgroundcolor=\color{black!2},xleftmargin=1em,framexleftmargin=0.5em,numbers=none,aboveskip=0.5em,belowskip=0.3em]
$( demo0.mm  1-Jan-04 $)

$(
                      PUBLIC DOMAIN DEDICATION

This file is placed in the public domain per the Creative Commons Public
Domain Dedication. http://creativecommons.org/licenses/publicdomain/

Norman Megill - email: nm at alum.mit.edu

$)


$( This file is the introductory formal system example described
   in Chapter 2 of the Metamath book. $)

$( Declare the constant symbols we will use $)
    $c 0 + = -> ( ) term wff |- $.
$( Declare the metavariables we will use $)
    $v t r s P Q $.
$( Specify properties of the metavariables $)
    tt $f term t $.
    tr $f term r $.
    ts $f term s $.
    wp $f wff P $.
    wq $f wff Q $.
$( Define "term" (part 1) $)
    tze $a term 0 $.
$( Define "term" (part 2) $)
    tpl $a term ( t + r ) $.
$( Define "wff" (part 1) $)
    weq $a wff t = r $.
$( Define "wff" (part 2) $)
    wim $a wff ( P -> Q ) $.
$( State axiom a1 $)
    a1 $a |- ( t = r -> ( t = s -> r = s ) ) $.
$( State axiom a2 $)
    a2 $a |- ( t + 0 ) = t $.
    ${
       min $e |- P $.
       maj $e |- ( P -> Q ) $.
$( Define the modus ponens inference rule $)
       mp  $a |- Q $.
    $}
$( Prove a theorem $)
    th1 $p |- t = t $=
  $( Here is its proof: $)
       tt tze tt tpl weq tt tt weq tt a2 tt tze tpl
       tt weq tt tze tpl tt weq tt tt weq wim tt a2
       tt tze tpl tt tt a1 mp mp
     $.
\end{lstlisting}

\noindent\textbf{issue107.mm} [REJECT] \textit{4.2.1: proof verification failure (missing hypothesis)}\par

\begin{lstlisting}[language=Metamath,basicstyle=\ttfamily\scriptsize,frame=single,framerule=0.3pt,rulecolor=\color{black!20},backgroundcolor=\color{black!2},xleftmargin=1em,framexleftmargin=0.5em,numbers=none,aboveskip=0.5em,belowskip=0.3em]
$c axiom hypothesis thesis $.

ax $a axiom $.

${
   hyp $e hypothesis $.
   mp $a thesis $.
$}

th $p thesis $= ax mp $.
\end{lstlisting}

\subsection{Active Hypotheses (\S4.2.5)}

\noindent\textbf{test23\_using\_non-essential\_hypotheses.mm} [REJECT] \textit{4.2.5: only active hypotheses may be used}\par

\begin{lstlisting}[language=Metamath,basicstyle=\ttfamily\scriptsize,frame=single,framerule=0.3pt,rulecolor=\color{black!20},backgroundcolor=\color{black!2},xleftmargin=1em,framexleftmargin=0.5em,numbers=none,aboveskip=0.5em,belowskip=0.3em]
$( Unit Test 23: Using non-essential hypotheses $)
$( Should reject: True $)

$c wff |- $.
$v x $.
wf $f wff x $.
${
  hyp $e |- x $.
  th1 $p |- x $= hyp $.
$}
evil $p |- x $= ( hyp ) A $.
\end{lstlisting}

\subsection{\$d Symbol Requirements (\S4.2.6)}

\noindent\textbf{test69\_djvars\_error\_masking.mm} [REJECT] \textit{4.2.6: \textbackslash{}\$d symbols must be declared variables}\par

\begin{lstlisting}[language=Metamath,basicstyle=\ttfamily\scriptsize,frame=single,framerule=0.3pt,rulecolor=\color{black!20},backgroundcolor=\color{black!2},xleftmargin=1em,framexleftmargin=0.5em,numbers=none,aboveskip=0.5em,belowskip=0.3em]
$( Test 69: $d with undeclared symbol should reject. $)
$( This currently exposes a parser diagnostic masking path in mm-lean4:
   the immediate token error can be overwritten by EOF unclosed-$d reporting. $)

$c wff $.
$v x $.

$d x y $.
\end{lstlisting}

\subsection{Disjoint Variable Violations (\S4.2.7)}

\noindent\textbf{disjoint1.mm} [REJECT] \textit{4.2.7: \textbackslash{}\$d violation in proof}\par

\begin{lstlisting}[language=Metamath,basicstyle=\ttfamily\scriptsize,frame=single,framerule=0.3pt,rulecolor=\color{black!20},backgroundcolor=\color{black!2},xleftmargin=1em,framexleftmargin=0.5em,numbers=none,aboveskip=0.5em,belowskip=0.3em]
$c var formula |- $.
$v x y $.
xf $f formula x $.
yf $f formula y $.
${
  $d x y $.
  ax-1 $a |- x y $.
$}
bad $p |- x x $=
  xf xf ax-1 $.
\end{lstlisting}

\noindent\textbf{disjoint2.mm} [REJECT] \textit{4.2.7: \textbackslash{}\$d violation in proof}\par

\begin{lstlisting}[language=Metamath,basicstyle=\ttfamily\scriptsize,frame=single,framerule=0.3pt,rulecolor=\color{black!20},backgroundcolor=\color{black!2},xleftmargin=1em,framexleftmargin=0.5em,numbers=none,aboveskip=0.5em,belowskip=0.3em]
$c formula |- ( ) $.

$v x y z w $.
xf $f formula x $.
yf $f formula y $.
zf $f formula z $.
wf $f formula w $.
combo $a formula ( x y ) $.
${
  $d x y $.
  ax-1 $a |- ( x y ) $.
$}

verybad $p |- ( ( x y ) ( z x ) ) $=
  xf yf combo zf xf combo ax-1 $.

${
  $d x y z $.
  stillbad $p |- ( ( x y ) ( z x ) ) $=
    xf yf combo zf xf combo ax-1 $.
$}

semigood $p |- ( ( x y ) ( z w ) ) $=
  xf yf combo zf wf combo ax-1 $.

${
  $d x y $.
  $d z w $.
  almostgood $p |- ( ( x y ) ( z w ) ) $=
    xf yf combo zf wf combo ax-1 $.
$}

${
  $d x y z w $.
  good $p |- ( ( x y ) ( z w ) ) $=
    xf yf combo zf wf combo ax-1 $.
$}

${
  $d x z $. $d x w $. $d y z $. $d y w $.
  stillgood $p |- ( ( x y ) ( z w ) ) $=
    xf yf combo zf wf combo ax-1 $.
$}
\end{lstlisting}

\noindent\textbf{disjoint3.mm} [REJECT] \textit{4.2.7: \textbackslash{}\$d violation in proof}\par

\begin{lstlisting}[language=Metamath,basicstyle=\ttfamily\scriptsize,frame=single,framerule=0.3pt,rulecolor=\color{black!20},backgroundcolor=\color{black!2},xleftmargin=1em,framexleftmargin=0.5em,numbers=none,aboveskip=0.5em,belowskip=0.3em]
$c var formula |- $.
$v x y $.
xf $f formula x $.
yf $f formula y $.
${
  $d x y $.
  ax-1 $a |- x y $.
$}
bad $p |- x y $=
  xf yf ax-1 $.
\end{lstlisting}

\noindent\textbf{test27\_disjoint\_variable\_constraint\_violation.mm} [REJECT] \textit{4.2.7: \textbackslash{}\$d constraints must be satisfied}\par

\begin{lstlisting}[language=Metamath,basicstyle=\ttfamily\scriptsize,frame=single,framerule=0.3pt,rulecolor=\color{black!20},backgroundcolor=\color{black!2},xleftmargin=1em,framexleftmargin=0.5em,numbers=none,aboveskip=0.5em,belowskip=0.3em]
$( Unit Test 27: Disjoint variable constraint violation $)
$( Should reject: True $)

$c wff -> |- ( ) $.
$v x y z $.
wfx $f wff x $.
wfy $f wff y $.
wfz $f wff z $.
$d x y $.

$( Axiom mentions BOTH x and y, so $d x y is mandatory for it $)
axxy $a |- ( x -> y ) $.

$( This proof violates $d: substitutes z for both x and y $)
$( x:=z, y:=z violates $d x y $)
bad $p |- ( z -> z ) $= wfz wfz axxy $.
\end{lstlisting}

\subsection{Include Processing (\S4.2.8)}

\noindent\textbf{test17\_include\_scope\_violation.mm} [REJECT] \textit{4.2.8: \textbackslash{}\$[ \textbackslash{}\$] scope rules}\par

\begin{lstlisting}[language=Metamath,basicstyle=\ttfamily\scriptsize,frame=single,framerule=0.3pt,rulecolor=\color{black!20},backgroundcolor=\color{black!2},xleftmargin=1em,framexleftmargin=0.5em,numbers=none,aboveskip=0.5em,belowskip=0.3em]
$( Unit Test 17: Include inside block with scope violation $)
$( Should reject: True $)

$c wff |- $.
$v x $.
wx $f wff x $.

${
  $( Include inside block - contents scoped to block $)
  $[ ./helpers/test17_helper.mm $]

  $( This would work - using inside the block $)
  $( th1 $p |- y $= wy ax-inner $. $)
$}

$( Try to use ax-inner outside the block - SCOPE VIOLATION $)
th2 $p |- y $= wy ax-inner $.
\end{lstlisting}

\noindent\textbf{test28\_self\_include.mm} [REJECT] \textit{4.2.8: file may not include itself}\par

\begin{lstlisting}[language=Metamath,basicstyle=\ttfamily\scriptsize,frame=single,framerule=0.3pt,rulecolor=\color{black!20},backgroundcolor=\color{black!2},xleftmargin=1em,framexleftmargin=0.5em,numbers=none,aboveskip=0.5em,belowskip=0.3em]
$( Unit Test 28: Self-include $)
$( Should reject: False - Spec section 4.1.2 says "will simply be ignored" $)
$( Note: metamath.exe REJECTS (spec divergence), mmverify_pure.py ACCEPTS (spec-compliant) $)

$c wff $.

$( Include this file itself - causes duplicate declarations $)
$[ ./test28_self_include.mm $]
\end{lstlisting}

\noindent\textbf{test44\_include\_cycle\_main.mm} [REJECT] \textit{4.2.8: circular include detected}\par

\begin{lstlisting}[language=Metamath,basicstyle=\ttfamily\scriptsize,frame=single,framerule=0.3pt,rulecolor=\color{black!20},backgroundcolor=\color{black!2},xleftmargin=1em,framexleftmargin=0.5em,numbers=none,aboveskip=0.5em,belowskip=0.3em]
$( Unit Test 44: Cycle detection - A includes B, B includes A $)
$( Should reject: False - cycle should be detected and ignored $)
$( Spec: Include stack tracking prevents infinite loops $)

$c wff |- $.

$( Include A, which includes B, which tries to include A again $)
$[ ./helpers/test44_helper_a.mm $]

$( Both A and B should be included exactly once $)
$( Use declarations from both files $)
th-a $p |- var-a $= fa-cycle ax-a-cycle $.
th-b $p |- var-b $= fb-cycle ax-b-cycle $.
\end{lstlisting}

\noindent\textbf{test46\_duplicate\_include\_main.mm} [REJECT] \textit{4.2.8: duplicate include processing}\par

\begin{lstlisting}[language=Metamath,basicstyle=\ttfamily\scriptsize,frame=single,framerule=0.3pt,rulecolor=\color{black!20},backgroundcolor=\color{black!2},xleftmargin=1em,framexleftmargin=0.5em,numbers=none,aboveskip=0.5em,belowskip=0.3em]
$( Unit Test 46: Duplicate include in different inner blocks $)
$( Should reject: True - includes must be at outermost level $)
$( Category: CORE - Strict include semantics $)
$( Spec Section 4.1.2: Include directives processed at outermost level only $)

$c wff |- $.

${
  $( First include: processes the file $)
  $[ ./helpers/test46_helper.mm $]

  $( Use the included content $)
  th1 $p |- y $= wy ax-y $.
$}

${
  $( Second include: should be ignored as whitespace $)
  $[ ./helpers/test46_helper.mm $]

  $( This would fail if file processed twice (y redeclared) $)
  $( But since second include ignored, ax-y and wy are still available globally $)
  th2 $p |- y $= wy ax-y $.
$}
\end{lstlisting}

\noindent\textbf{test52\_include\_inside\_e.mm} [REJECT] \textit{4.2.8: \textbackslash{}\$[ not allowed inside \textbackslash{}\$e}\par

\begin{lstlisting}[language=Metamath,basicstyle=\ttfamily\scriptsize,frame=single,framerule=0.3pt,rulecolor=\color{black!20},backgroundcolor=\color{black!2},xleftmargin=1em,framexleftmargin=0.5em,numbers=none,aboveskip=0.5em,belowskip=0.3em]
$( Test 52: $[ ... $] inside an $e statement - must REJECT per Sec 4.1.2 (includes only at outermost level) $)

$c wff |- $.
$v ph $.

e1 $e $[ ./helpers/test40_include_scope_correct_inner.mm $] $.
\end{lstlisting}

\noindent\textbf{test53\_include\_inside\_d.mm} [REJECT] \textit{4.2.8: \textbackslash{}\$[ not allowed inside \textbackslash{}\$d}\par

\begin{lstlisting}[language=Metamath,basicstyle=\ttfamily\scriptsize,frame=single,framerule=0.3pt,rulecolor=\color{black!20},backgroundcolor=\color{black!2},xleftmargin=1em,framexleftmargin=0.5em,numbers=none,aboveskip=0.5em,belowskip=0.3em]
$( Test 53: $[ ... $] inside a $d statement - must REJECT per Sec 4.1.2 (includes only at outermost level) and Sec 4.2.4 (varlist only) $)

$c wff |- $.
$v x y $.

$d $[ ./helpers/test40_include_scope_correct_inner.mm $] x y $.
\end{lstlisting}

\subsection{Compressed Proof Format (\S4.4)}

\noindent\textbf{out-of-range-saved-step-bad1.mm} [REJECT] \textit{4.4: compressed proof Z reference out of range}\par

\begin{lstlisting}[language=Metamath,basicstyle=\ttfamily\scriptsize,frame=single,framerule=0.3pt,rulecolor=\color{black!20},backgroundcolor=\color{black!2},xleftmargin=1em,framexleftmargin=0.5em,numbers=none,aboveskip=0.5em,belowskip=0.3em]
$( out-of-range-saved-step-bad1.mm
   Incorrect compressed proof: An example generated by ChatGPT o3
   That references a proof index greater than the number of saved proof steps.
   Added because mmverify.py was missing a raise MMError and tests did not flag it. $)

$c |- T $.

ax     $a |- T $.
buggy  $p |- T $= ( ax ) A Z C $.
\end{lstlisting}

\subsection{Compressed Proof Encoding (Appendix B)}

\noindent\textbf{test19\_illegal\_characters\_in\_compressed\_proof.mm} [REJECT] \textit{AppB: compressed proof uses only A-Z and ? (strict)}\par

\begin{lstlisting}[language=Metamath,basicstyle=\ttfamily\scriptsize,frame=single,framerule=0.3pt,rulecolor=\color{black!20},backgroundcolor=\color{black!2},xleftmargin=1em,framexleftmargin=0.5em,numbers=none,aboveskip=0.5em,belowskip=0.3em]
$( Unit Test 19: Illegal characters in compressed proof $)
$( Should reject: True $)

$c wff |- $.
$v x $.
wf $f wff x $.
ax $a |- x $.
th $p |- x $= ( ax ) A0B $.
\end{lstlisting}

\noindent\textbf{test29\_compressed\_proof\_header\_mismatch.mm} [REJECT] \textit{AppB: compressed proof header must match mandatory hyps}\par

\begin{lstlisting}[language=Metamath,basicstyle=\ttfamily\scriptsize,frame=single,framerule=0.3pt,rulecolor=\color{black!20},backgroundcolor=\color{black!2},xleftmargin=1em,framexleftmargin=0.5em,numbers=none,aboveskip=0.5em,belowskip=0.3em]
$( Unit Test 29: Compressed proof header mismatch $)
$( Should reject: True $)

$c wff |- $.
$v x $.
wf $f wff x $.
ax1 $a |- x $.
ax2 $a |- x $.
$( Header lists ax1 but proof uses ax2! $)
bad $p |- x $= ( ax1 ) B $.
\end{lstlisting}

\noindent\textbf{test36\_whitespace\_in\_compressed\_proof\_bad.mm} [REJECT] \textit{AppB: no whitespace inside compressed proof string}\par

\begin{lstlisting}[language=Metamath,basicstyle=\ttfamily\scriptsize,frame=single,framerule=0.3pt,rulecolor=\color{black!20},backgroundcolor=\color{black!2},xleftmargin=1em,framexleftmargin=0.5em,numbers=none,aboveskip=0.5em,belowskip=0.3em]
$( Unit Test 36: Whitespace in compressed proof $)
$( Should reject: True - This database has invalid proof steps B and C $)
$( Original intent was to test whitespace handling, but ABC is invalid $)

$c wff |- $.
$v x $.
wf $f wff x $.
ax $a |- x $.
$( Compressed proof ABC: A=0 (wf), B=1 (ax), C=2 (out of bounds!) $)
th $p |- x $= ( ax ) A
  B
	C $.
\end{lstlisting}

\noindent\textbf{test38\_whitespace\_in\_valid\_compressed.mm} [ACCEPT] \textit{AppB: whitespace allowed in header, not proof string}\par

\begin{lstlisting}[language=Metamath,basicstyle=\ttfamily\scriptsize,frame=single,framerule=0.3pt,rulecolor=\color{black!20},backgroundcolor=\color{black!2},xleftmargin=1em,framexleftmargin=0.5em,numbers=none,aboveskip=0.5em,belowskip=0.3em]
$( Unit Test 38: Whitespace in valid compressed proof $)
$( Should reject: False - whitespace should be ignored per spec section 4.4.2 $)

$c wff |- $.
$v x $.
wf $f wff x $.
ax $a |- x $.

$( Compressed proof with whitespace: spaces, newlines, tabs between valid steps $)
$( Proof is just: push wf (A), then ax (B) $)
$( Label list: [wf=0, ax=1], so A B is valid $)
th $p |- x $= ( ax )   A
  B   $.
\end{lstlisting}

\noindent\textbf{test55\_compressed\_z\_oob.mm} [REJECT] \textit{AppB: Z reference out of bounds}\par

\begin{lstlisting}[language=Metamath,basicstyle=\ttfamily\scriptsize,frame=single,framerule=0.3pt,rulecolor=\color{black!20},backgroundcolor=\color{black!2},xleftmargin=1em,framexleftmargin=0.5em,numbers=none,aboveskip=0.5em,belowskip=0.3em]
$( Test 55: Compressed proof index out of bounds - must REJECT per Appendix B (label index must be within header list) $)

$c wff |- $.
$v ph $.
f1 $f wff ph $.
a1 $a |- ph $.

p1 $p |- ph $= ( a1 ) B $.
\end{lstlisting}

\noindent\textbf{test56\_compressed\_z\_simple\_bad.mm} [REJECT] \textit{AppB: compressed Z out of bounds}\par

\begin{lstlisting}[language=Metamath,basicstyle=\ttfamily\scriptsize,frame=single,framerule=0.3pt,rulecolor=\color{black!20},backgroundcolor=\color{black!2},xleftmargin=1em,framexleftmargin=0.5em,numbers=none,aboveskip=0.5em,belowskip=0.3em]
$( Minimal test of compressed proof with Z (save) marker $)

$c wff |- $.
$v ph ps $.
wph $f wff ph $.
wps $f wff ps $.

$( A simple axiom - implication introduction $)
ax-1 $a |- ( ph -> ( ps -> ph ) ) $.

$( Axiom using ax-1 as basis $)
wi $a wff ( ph -> ps ) $.

$( Prove something using Z - save and reuse $)
$( This proof:
   - A = index 0 = wph (mandatory f-hyp for ph)
   - B = index 1 = wps (mandatory f-hyp for ps)
   - C = index 2 = wi (explicit label)
   - D = index 3 = ax-1 (explicit label)
   So: ( wi ax-1 ) A B C Z D means:
   1. Push wph -> get "wff ph" on stack
   2. Push wps -> get "wff ps" on stack
   3. Apply wi -> get "wff ( ph -> ps )" on stack
   4. Z = save top (which is "wff ( ph -> ps )")
   5. Apply ax-1 with substitutions -> get "|- ( ph -> ( ps -> ph ) )"
   Actually, let me simplify...
$)

$( Simple proof without Z first to verify basic compressed proofs work $)
simple1 $p wff ( ph -> ps ) $= ( wi ) AB $.

$( Now with Z $)
$.
\end{lstlisting}

\subsection{Valid Databases (Positive Tests)}

\noindent\textbf{anatomy.mm} [ACCEPT] \textit{Valid database per spec}\par

\begin{lstlisting}[language=Metamath,basicstyle=\ttfamily\scriptsize,frame=single,framerule=0.3pt,rulecolor=\color{black!20},backgroundcolor=\color{black!2},xleftmargin=1em,framexleftmargin=0.5em,numbers=none,aboveskip=0.5em,belowskip=0.3em]
$( anatomy.mm $)

$(
                      PUBLIC DOMAIN DEDICATION

This file is placed in the public domain per the Creative Commons Public
Domain Dedication. http://creativecommons.org/licenses/publicdomain/

Norman Megill - email: nm at alum.mit.edu

$)


$( This file is the introductory formal system example described
   in Chapter 4 of the Metamath book, "Anatomy of a proof". $)

$( Declare the constant symbols we will use $)
$c ( ) -> wff $.
$( Declare the metavariables we will use $)
$v p q r s $.
$( Declare wffs $)
wp $f wff p $.
wq $f wff q $.
wr $f wff r $.
ws $f wff s $.
w2 $a wff ( p -> q ) $.

$( Prove a simple theorem. $)
wnew $p wff ( s -> ( r -> p ) ) $= ws wr wp w2 w2 $.
\end{lstlisting}

\noindent\textbf{demo0-includee.mm} [ACCEPT] \textit{Valid database per spec}\par

\begin{lstlisting}[language=Metamath,basicstyle=\ttfamily\scriptsize,frame=single,framerule=0.3pt,rulecolor=\color{black!20},backgroundcolor=\color{black!2},xleftmargin=1em,framexleftmargin=0.5em,numbers=none,aboveskip=0.5em,belowskip=0.3em]
$( demo0-includee.mm  1-Jan-04 $)

$(
                      PUBLIC DOMAIN DEDICATION

This file is placed in the public domain per the Creative Commons Public
Domain Dedication. http://creativecommons.org/licenses/publicdomain/

Norman Megill - email: nm at alum.mit.edu

$)


$( This file is the introductory formal system example described
   in Chapter 2 of the Meamath book. $)

$( Declare the constant symbols we will use $)
    $c 0 + = -> ( ) term wff |- $.
$( Declare the metavariables we will use $)
    $v t r s P Q $.
$( Specify properties of the metavariables $)
    tt $f term t $.
    tr $f term r $.
    ts $f term s $.
    wp $f wff P $.
    wq $f wff Q $.
$( Define "term" (part 1) $)
    tze $a term 0 $.
$( Define "term" (part 2) $)
    tpl $a term ( t + r ) $.
$( Define "wff" (part 1) $)
    weq $a wff t = r $.
$( Define "wff" (part 2) $)
    wim $a wff ( P -> Q ) $.
$( State axiom a1 $)
    a1 $a |- ( t = r -> ( t = s -> r = s ) ) $.
$( State axiom a2 $)
    a2 $a |- ( t + 0 ) = t $.
    ${
       min $e |- P $.
       maj $e |- ( P -> Q ) $.
$( Define the modus ponens inference rule $)
       mp  $a |- Q $.
    $}
\end{lstlisting}

\noindent\textbf{demo0-includer.mm} [ACCEPT] \textit{Valid database per spec}\par

\begin{lstlisting}[language=Metamath,basicstyle=\ttfamily\scriptsize,frame=single,framerule=0.3pt,rulecolor=\color{black!20},backgroundcolor=\color{black!2},xleftmargin=1em,framexleftmargin=0.5em,numbers=none,aboveskip=0.5em,belowskip=0.3em]
$( demo0-includer.mm  1-Jan-04 $)

$(
                      PUBLIC DOMAIN DEDICATION

This file is placed in the public domain per the Creative Commons Public
Domain Dedication. http://creativecommons.org/licenses/publicdomain/

Norman Megill - email: nm at alum.mit.edu

$)


$[ demo0-includee.mm $]

$( Prove a theorem $)
    th1 $p |- t = t $=
  $( Here is its proof: $)
       tt tze tpl tt weq tt tt weq tt a2 tt tze tpl
       tt weq tt tze tpl tt weq tt tt weq wim tt a2
       tt tze tpl tt tt a1 mp mp
     $.
\end{lstlisting}

\noindent\textbf{demo0.mm} [ACCEPT] \textit{Valid database per spec}\par

\begin{lstlisting}[language=Metamath,basicstyle=\ttfamily\scriptsize,frame=single,framerule=0.3pt,rulecolor=\color{black!20},backgroundcolor=\color{black!2},xleftmargin=1em,framexleftmargin=0.5em,numbers=none,aboveskip=0.5em,belowskip=0.3em]
$( demo0.mm  1-Jan-04 $)

$(
                      PUBLIC DOMAIN DEDICATION

This file is placed in the public domain per the Creative Commons Public
Domain Dedication. http://creativecommons.org/licenses/publicdomain/

Norman Megill - email: nm at alum.mit.edu

$)


$( This file is the introductory formal system example described
   in Chapter 2 of the Meamath book. $)

$( Declare the constant symbols we will use $)
    $c 0 + = -> ( ) term wff |- $.
$( Declare the metavariables we will use $)
    $v t r s P Q $.
$( Specify properties of the metavariables $)
    tt $f term t $.
    tr $f term r $.
    ts $f term s $.
    wp $f wff P $.
    wq $f wff Q $.
$( Define "term" (part 1) $)
    tze $a term 0 $.
$( Define "term" (part 2) $)
    tpl $a term ( t + r ) $.
$( Define "wff" (part 1) $)
    weq $a wff t = r $.
$( Define "wff" (part 2) $)
    wim $a wff ( P -> Q ) $.
$( State axiom a1 $)
    a1 $a |- ( t = r -> ( t = s -> r = s ) ) $.
$( State axiom a2 $)
    a2 $a |- ( t + 0 ) = t $.
    ${
       min $e |- P $.
       maj $e |- ( P -> Q ) $.
$( Define the modus ponens inference rule $)
       mp  $a |- Q $.
    $}
$( Prove a theorem $)
    th1 $p |- t = t $=
  $( Here is its proof: $)
       tt tze tpl tt weq tt tt weq tt a2 tt tze tpl
       tt weq tt tze tpl tt weq tt tt weq wim tt a2
       tt tze tpl tt tt a1 mp mp
     $.
\end{lstlisting}

\noindent\textbf{empty.mm} [ACCEPT] \textit{Empty database is valid}\par

\begin{lstlisting}[language=Metamath,basicstyle=\ttfamily\scriptsize,frame=single,framerule=0.3pt,rulecolor=\color{black!20},backgroundcolor=\color{black!2},xleftmargin=1em,framexleftmargin=0.5em,numbers=none,aboveskip=0.5em,belowskip=0.3em]

\end{lstlisting}

\noindent\textbf{issue129.mm} [ACCEPT] \textit{Valid database per spec}\par

\begin{lstlisting}[language=Metamath,basicstyle=\ttfamily\scriptsize,frame=single,framerule=0.3pt,rulecolor=\color{black!20},backgroundcolor=\color{black!2},xleftmargin=1em,framexleftmargin=0.5em,numbers=none,aboveskip=0.5em,belowskip=0.3em]
$c A $.

ax-1 $a A A A A A $.

th1 $p A A A A A $= ax-1 $.

$( $t
    latexdef "A" as "\text{ A A A A A A A A A A }";
$)
\end{lstlisting}

\noindent\textbf{issue134.mm} [ACCEPT] \textit{Valid database per spec}\par

\begin{lstlisting}[language=Metamath,basicstyle=\ttfamily\scriptsize,frame=single,framerule=0.3pt,rulecolor=\color{black!20},backgroundcolor=\color{black!2},xleftmargin=1em,framexleftmargin=0.5em,numbers=none,aboveskip=0.5em,belowskip=0.3em]
$c A $.
ax1 $a A $.
\end{lstlisting}

\noindent\textbf{local-var-good.mm} [ACCEPT] \textit{Valid variable scope per spec}\par

\begin{lstlisting}[language=Metamath,basicstyle=\ttfamily\scriptsize,frame=single,framerule=0.3pt,rulecolor=\color{black!20},backgroundcolor=\color{black!2},xleftmargin=1em,framexleftmargin=0.5em,numbers=none,aboveskip=0.5em,belowskip=0.3em]
  $c formula $.
  $c A. $.
  $c setvar $.
  $v ph $.
  wph $f formula ph $.

  ${
    $v x $.
    $( Let ` x ` be an individual variable (temporary declaration). $)
    vx.wal $f setvar x $.
    $( Extend formula definition to include the universal quantifier ('for all').  $)
    wal $a formula A. x ph $.
  $}
\end{lstlisting}

\noindent\textbf{min\_found.mm} [ACCEPT] \textit{Valid database per spec}\par

\begin{lstlisting}[language=Metamath,basicstyle=\ttfamily\scriptsize,frame=single,framerule=0.3pt,rulecolor=\color{black!20},backgroundcolor=\color{black!2},xleftmargin=1em,framexleftmargin=0.5em,numbers=none,aboveskip=0.5em,belowskip=0.3em]
$( A shortening found by the minimizer.  The comparison with min_not_found.mm
   suggests that the minimizer's ability to detect it depends on the presence
   of statement "A" in the steps of the minimized proof.  The costant C and
   axiom ax-3 are not used here, but they are kept anyway to lessen the amount
   of differences with min_not_found.mm. $)

  $c A B C D $.

  ax-1 $a A $.
  ax-2 $a B $.
  ax-3 $a C $.

  ${
    in1.1 $e A $.
    in1 $a D $.
  $}

  ${
    in2.1 $e B $.
    in2.2 $e A $.
    in2.3 $e B $.
    in2.4 $e B $.
    in2.5 $e B $.
    in2 $a D $.
  $}

  $( Minimized proof, found automatically.  (New usage is discouraged.) $)
  th1 $p D $= ( ax-1 in1 ) AB $.

  $( Non-minimized proof. $)
  th2 $p D $= ( ax-2 ax-1 in2 ) ABAAAC $.
\end{lstlisting}

\noindent\textbf{min\_not\_found.mm} [ACCEPT] \textit{Valid database per spec}\par

\begin{lstlisting}[language=Metamath,basicstyle=\ttfamily\scriptsize,frame=single,framerule=0.3pt,rulecolor=\color{black!20},backgroundcolor=\color{black!2},xleftmargin=1em,framexleftmargin=0.5em,numbers=none,aboveskip=0.5em,belowskip=0.3em]
$( A shortening not found by the minimizer.  The comparison with min_found.mm
   suggests that the minimizer's inability to detect it is due to the absence
   of statement "C" in any of the steps of the minimized proof.  The
   shortening is detected if proveFloating is set to 1 in replaceStatement
   called by minimizeProof, which allows absent $e statements to be found. $)

  $c A B C D $.

  ax-1 $a A $.
  ax-2 $a B $.
  ax-3 $a C $.

  ${
    in1.1 $e A $.
    in1 $a D $.
  $}

  ${
    in2.1 $e B $.
    in2.2 $e C $.
    in2.3 $e B $.
    in2.4 $e B $.
    in2.5 $e B $.
    in2 $a D $.
  $}

  $( Minimized proof, not found automatically.  (New usage is discouraged.) $)
  th1 $p D $= ( ax-1 in1 ) AB $.

  $( Non-minimized proof. $)
  th2 $p D $= ( ax-2 ax-3 in2 ) ABAAAC $.
\end{lstlisting}

\noindent\textbf{nuniq-gram.mm} [ACCEPT] \textit{Valid database per spec (non-unique grammar allowed)}\par

\begin{lstlisting}[language=Metamath,basicstyle=\ttfamily\scriptsize,frame=single,framerule=0.3pt,rulecolor=\color{black!20},backgroundcolor=\color{black!2},xleftmargin=1em,framexleftmargin=0.5em,numbers=none,aboveskip=0.5em,belowskip=0.3em]
$( Here we have two variable typecodes A and B, and one logical
   typecode |-, and the question is whether |- connects to A or B.
   It parses either way. This is actually an unambiguous formal system,
   because we require only that there are no two ways to derive the proof
   of the *same* syntax theorem, and "A *" and "B *" are different theorems. $)
$c A B |- * $. $v a b $.
va $f A a $.
vb $f B b $.
as $a A * $.
bs $a B * $.
ax $a |- * $.
\end{lstlisting}

\noindent\textbf{test37\_rpn\_interleaving\_mandatory\_hyps.mm} [ACCEPT] \textit{Valid per RPN proof semantics}\par

\begin{lstlisting}[language=Metamath,basicstyle=\ttfamily\scriptsize,frame=single,framerule=0.3pt,rulecolor=\color{black!20},backgroundcolor=\color{black!2},xleftmargin=1em,framexleftmargin=0.5em,numbers=none,aboveskip=0.5em,belowskip=0.3em]
$( Unit Test 37: RPN interleaving of $f and $e in mandatory hypotheses $)
$( Should reject: False - This documents correct behavior $)

$c wff |- -> ( ) $.
$v ph ps $.

${
  $( $f for ph $)
  wph $f wff ph $.

  $( $e using ph $)
  h1 $e |- ph $.

  $( $f for ps - AFTER the $e (interleaved!) $)
  wps $f wff ps $.

  $( Assertion: mandatory hyps in RPN order MUST be: wph, h1, wps $)
  $( NOT: wph, wps, h1 (which would be "all $f then $e") $)
  ax-test $a |- ( ph -> ps ) $.
$}

$( This test documents the CORRECT mandatory hypothesis order. $)
$( A verifier that incorrectly orders as "all $f, then all $e" would get: $)
$(   WRONG: wph, wps, h1 $)
$( The correct order per spec is appearance order: $)
$(   CORRECT: wph, h1, wps $)

$( To verify: use metamath.exe "show statement ax-test /full" $)
$( Expected output: "Its mandatory hypotheses in RPN order are:" $)
$(   wph $f wff ph $. $)
$(   h1 $e |- ph $. $)
$(   wps $f wff ps $. $)
\end{lstlisting}

\noindent\textbf{test39\_include\_basic\_outer.mm} [ACCEPT] \textit{Valid include at outermost scope}\par

\begin{lstlisting}[language=Metamath,basicstyle=\ttfamily\scriptsize,frame=single,framerule=0.3pt,rulecolor=\color{black!20},backgroundcolor=\color{black!2},xleftmargin=1em,framexleftmargin=0.5em,numbers=none,aboveskip=0.5em,belowskip=0.3em]
$( Unit Test 39: Basic include - outer file $)
$( Should reject: False - legitimate include should work $)

$( Include provides constants, variables, and axioms $)
$[ ./helpers/test39_helper.mm $]

$( Use the included axiom to prove something $)
th1 $p |- ph $= wph ax-1 $.
\end{lstlisting}

\noindent\textbf{test\_complex\_ehyp\_syl2anc.mm} [ACCEPT] \textit{Valid proof per spec}\par

\begin{lstlisting}[language=Metamath,basicstyle=\ttfamily\scriptsize,frame=single,framerule=0.3pt,rulecolor=\color{black!20},backgroundcolor=\color{black!2},xleftmargin=1em,framexleftmargin=0.5em,numbers=none,aboveskip=0.5em,belowskip=0.3em]
$( Minimized test case for syl2anc - tests complex essential hypothesis handling $)
$( Based on hol.mm by Mario Carneiro $)
$( This proof SHOULD verify correctly $)

$c term |- |= ( ) , $.
$v R S T A B $.

tr $f term R $.
ts $f term S $.
tt $f term T $.
ta $f term A $.
tb $f term B $.

$( Syntax: context pair $)
kct $a term ( A , B ) $.

$( Axioms with essential hypotheses $)
${
  ax-syl.1 $e |- R |= S $.
  ax-syl.2 $e |- S |= T $.
  ax-syl $a |- R |= T $.
$}

${
  ax-jca.1 $e |- R |= S $.
  ax-jca.2 $e |- R |= T $.
  ax-jca $a |- R |= ( S , T ) $.
$}

$( Derived theorems $)
${
  syl.1 $e |- R |= S $.
  syl.2 $e |- S |= T $.
  syl $p |- R |= T $=
    tr ts tt syl.1 syl.2 ax-syl $.
$}

${
  jca.1 $e |- R |= S $.
  jca.2 $e |- R |= T $.
  jca $p |- R |= ( S , T ) $=
    tr ts tt jca.1 jca.2 ax-jca $.
$}

${
  syl2anc.1 $e |- R |= S $.
  syl2anc.2 $e |- R |= T $.
  syl2anc.3 $e |- ( S , T ) |= A $.
  syl2anc $p |- R |= A $=
    tr ts tt kct ta tr ts tt syl2anc.1 syl2anc.2 jca syl2anc.3 syl $.
$}
\end{lstlisting}

\noindent\textbf{test59\_f\_type\_global\_conflict.mm} [ACCEPT] \textit{SPEC DIVERGENCE: L156-158 says REJECT, all impls ACCEPT}\par

\begin{lstlisting}[language=Metamath,basicstyle=\ttfamily\scriptsize,frame=single,framerule=0.3pt,rulecolor=\color{black!20},backgroundcolor=\color{black!2},xleftmargin=1em,framexleftmargin=0.5em,numbers=none,aboveskip=0.5em,belowskip=0.3em]
$( test59_f_type_global_conflict.mm - ACCEPT (SPEC DIVERGENCE!)
   Spec 4.1.3 L156-158 says:
   "The type declared by a $f statement for a given label is global even if
   the variable is not (e.g., a database may not have wff P in one local
   scope and class P in another)."

   HOWEVER: All three verifiers (metamath.exe, metamath-knife, mmverify.py)
   ACCEPT having the same variable with different types in different scopes!

   This is documented as a SPEC DIVERGENCE: the spec says REJECT, but
   all implementations ACCEPT. The convention in set.mm is to never do this,
   so it's never been caught.

   For our test suite, we follow the IMPLEMENTATIONS (ACCEPT), not the spec.
$)

$c wff class |- $.

$( First scope: x is wff $)
${
  $v x $.
  wx $f wff x $.
  ax1 $a wff x $.
$}

$( Second scope: x is class - spec says REJECT, implementations say ACCEPT $)
${
  $v x $.
  cx $f class x $.
  ax2 $a class x $.
$}
\end{lstlisting}

\noindent\textbf{test60\_empty\_block.mm} [ACCEPT] \textit{EBNF: empty block \textbackslash{}\$\{ \textbackslash{}\$\} is valid}\par

\begin{lstlisting}[language=Metamath,basicstyle=\ttfamily\scriptsize,frame=single,framerule=0.3pt,rulecolor=\color{black!20},backgroundcolor=\color{black!2},xleftmargin=1em,framexleftmargin=0.5em,numbers=none,aboveskip=0.5em,belowskip=0.3em]
$( test60_empty_block.mm - ACCEPT
   Tests EBNF: block ::= '${' stmt* '$}'
   Zero statements in a block is valid.

   Empty blocks between valid proofs should be accepted.
$)

$c wff |- $.
$v p q $.
wp $f wff p $.
wq $f wff q $.

$( First proof $)
th1 $a wff p $.

$( Empty block - should be valid $)
${ $}

$( Another proof after empty block $)
th2 $a wff q $.

$( Nested empty blocks $)
${
  ${ $}
  ${ $}
$}

$( Final proof $)
th3 $p wff p $= wp $.
\end{lstlisting}

\noindent\textbf{test57\_compressed\_z\_minimal.mm} [ACCEPT] \textit{Valid compressed proof with Z}\par

\begin{lstlisting}[language=Metamath,basicstyle=\ttfamily\scriptsize,frame=single,framerule=0.3pt,rulecolor=\color{black!20},backgroundcolor=\color{black!2},xleftmargin=1em,framexleftmargin=0.5em,numbers=none,aboveskip=0.5em,belowskip=0.3em]
$( Minimal test of compressed proof with Z (save) marker $)
$( This is designed to be the simplest possible Z test that should pass $)

$c wff ( ) -> $.
$v ph ps $.
wph $f wff ph $.
wps $f wff ps $.

$( Binary connective - implication $)
wi $a wff ( ph -> ps ) $.

$( Simple proof using Z:
   Prove: wff ( ( ph -> ps ) -> ( ph -> ps ) )

   Labels: wph=0, wps=1, wi=2

   Proof steps:
   A = 0 = push wph -> stack: [wff ph]
   B = 1 = push wps -> stack: [wff ph, wff ps]
   C = 2 = apply wi -> stack: [wff ( ph -> ps )]
   Z = -1 = save top -> saved: [wff ( ph -> ps )], stack: [wff ( ph -> ps )]
   D = 3 = recall saved[0] -> stack: [wff ( ph -> ps ), wff ( ph -> ps )]
   C = 2 = apply wi -> stack: [wff ( ( ph -> ps ) -> ( ph -> ps ) )]

   Compressed: ABCZDC
$)
th1 $p wff ( ( ph -> ps ) -> ( ph -> ps ) ) $= ( wi ) ABCZDC $.
\end{lstlisting}

\noindent\textbf{test\_compressed\_simple.mm} [ACCEPT] \textit{Valid compressed proof}\par

\begin{lstlisting}[language=Metamath,basicstyle=\ttfamily\scriptsize,frame=single,framerule=0.3pt,rulecolor=\color{black!20},backgroundcolor=\color{black!2},xleftmargin=1em,framexleftmargin=0.5em,numbers=none,aboveskip=0.5em,belowskip=0.3em]
$( Minimal test for compressed proof $)

$c term |- |= $.
$v R S T $.

tR $f term R $.
tS $f term S $.
tT $f term T $.

${
  syl.1 $e |- R |= S $.
  syl.2 $e |- S |= T $.
  ax-syl $a |- R |= T $.
$}

${
  th.1 $e |- R |= S $.
  th.2 $e |- S |= T $.
  $( Uncompressed proof - should work $)
  th $p |- R |= T $= tR tS tT th.1 th.2 ax-syl $.
$}
\end{lstlisting}

\noindent\textbf{test\_dv\_pairwise.mm} [ACCEPT] \textit{Valid pairwise \textbackslash{}\$d}\par

\begin{lstlisting}[language=Metamath,basicstyle=\ttfamily\scriptsize,frame=single,framerule=0.3pt,rulecolor=\color{black!20},backgroundcolor=\color{black!2},xleftmargin=1em,framexleftmargin=0.5em,numbers=none,aboveskip=0.5em,belowskip=0.3em]
$( Test for DV pairwise expansion: $d x y z $. should create pairs (x,y), (x,z), (y,z) $)

$c wff |- $.
$v x y z $.

wx $f wff x $.
wy $f wff y $.
wz $f wff z $.

${
  $d x y z $.

  $( Axiom requiring x,y disjoint $)
  ax-xy.1 $e |- x $.
  ax-xy $a |- y $.
$}

${
  $d x y z $.

  $( Axiom requiring x,z disjoint $)
  ax-xz.1 $e |- x $.
  ax-xz $a |- z $.
$}

${
  $d x y z $.

  $( Axiom requiring y,z disjoint - THIS IS THE KEY TEST $)
  ax-yz.1 $e |- y $.
  ax-yz $a |- z $.
$}

$( Test theorem - uses all three DV pairs $)
${
  $d x y z $.
  th.1 $e |- x $.
  th $p |- z $= wx wz th.1 ax-xz $.
$}
\end{lstlisting}

\noindent\textbf{test\_dv\_yz\_required.mm} [ACCEPT] \textit{Valid \textbackslash{}\$d test}\par

\begin{lstlisting}[language=Metamath,basicstyle=\ttfamily\scriptsize,frame=single,framerule=0.3pt,rulecolor=\color{black!20},backgroundcolor=\color{black!2},xleftmargin=1em,framexleftmargin=0.5em,numbers=none,aboveskip=0.5em,belowskip=0.3em]
$( Test for DV pairwise expansion: $d x y z $. should create pair (y,z)
   This test REQUIRES (y,z) to be disjoint - will FAIL if pairwise expansion broken $)

$c wff |- $.
$v x y z $.

wx $f wff x $.
wy $f wff y $.
wz $f wff z $.

$( Key axiom: requires y and z to be disjoint $)
${
  $d y z $.
  ax-yz.1 $e |- y $.
  ax-yz $a |- z $.
$}

$( Theorem using multi-variable $d declaration
   The $d x y z $. MUST expand to include (y,z) for this to verify! $)
${
  $d x y z $.
  th.1 $e |- y $.
  th $p |- z $= wy wz th.1 ax-yz $.
$}
\end{lstlisting}

\noindent\textbf{test\_stack\_fhyps.mm} [ACCEPT] \textit{Valid stack test}\par

\begin{lstlisting}[language=Metamath,basicstyle=\ttfamily\scriptsize,frame=single,framerule=0.3pt,rulecolor=\color{black!20},backgroundcolor=\color{black!2},xleftmargin=1em,framexleftmargin=0.5em,numbers=none,aboveskip=0.5em,belowskip=0.3em]
$( Test for floating hypothesis consumption from stack $)

$c wff |- term |= $.
$v R S T $.

wR $f term R $.
wS $f term S $.
wT $f term T $.

${
  ax-syl.1 $e |- R |= S $.
  ax-syl.2 $e |- S |= T $.
  $( Syllogism: if R implies S and S implies T, then R implies T $)
  ax-syl $a |- R |= T $.
$}

${
  syl.1 $e |- R |= S $.
  syl.2 $e |- S |= T $.
  $( Proof using ax-syl - should need 5 items on stack:
     3 floating hyps (term R, term S, term T) + 2 essential hyps $)
  syl $p |- R |= T $= wR wS wT syl.1 syl.2 ax-syl $.
$}
\end{lstlisting}

\noindent\textbf{test\_stack\_simple.mm} [ACCEPT] \textit{Valid stack test}\par

\begin{lstlisting}[language=Metamath,basicstyle=\ttfamily\scriptsize,frame=single,framerule=0.3pt,rulecolor=\color{black!20},backgroundcolor=\color{black!2},xleftmargin=1em,framexleftmargin=0.5em,numbers=none,aboveskip=0.5em,belowskip=0.3em]
$( Simple test for stack operations $)

$c wff |- $.
$v x y $.

wx $f wff x $.
wy $f wff y $.

${
  ax $a |- x $.
  th $p |- y $= wy ax $.
$}
\end{lstlisting}

\subsection{Miscellaneous}

\noindent\textbf{test41\_include\_multiple\_main.mm} [ACCEPT] \textit{Multiple valid includes at outermost scope}\par

\begin{lstlisting}[language=Metamath,basicstyle=\ttfamily\scriptsize,frame=single,framerule=0.3pt,rulecolor=\color{black!20},backgroundcolor=\color{black!2},xleftmargin=1em,framexleftmargin=0.5em,numbers=none,aboveskip=0.5em,belowskip=0.3em]
$( Unit Test 41: Multiple includes $)
$( Should reject: False - multiple includes should work $)

$c wff |- $.

$( Include file A $)
$[ ./helpers/test41_helper_a.mm $]

$( Include file B $)
$[ ./helpers/test41_helper_b.mm $]

$( Use content from both includes $)
th-a $p |- var-a $= fa ax-a $.
th-b $p |- var-b $= fb ax-b $.
\end{lstlisting}

\noindent\textbf{test70\_djvars\_before\_float.mm} [ACCEPT] \textit{Metamath accepts \textbackslash{}\$d before corresponding \textbackslash{}\$f declarations}\par

\begin{lstlisting}[language=Metamath,basicstyle=\ttfamily\scriptsize,frame=single,framerule=0.3pt,rulecolor=\color{black!20},backgroundcolor=\color{black!2},xleftmargin=1em,framexleftmargin=0.5em,numbers=none,aboveskip=0.5em,belowskip=0.3em]
$( Test 70: $d may appear before corresponding $f hypotheses. $)

$c wff |- $.
$v x y $.
$d x y $.

hx $f wff x $.
hy $f wff y $.

a1 $a |- x $.
\end{lstlisting}

\noindent\textbf{test58\_compressed\_ehyp\_z.mm} [ACCEPT] \textit{compressed proof with essential hyps and Z markers}\par

\begin{lstlisting}[language=Metamath,basicstyle=\ttfamily\scriptsize,frame=single,framerule=0.3pt,rulecolor=\color{black!20},backgroundcolor=\color{black!2},xleftmargin=1em,framexleftmargin=0.5em,numbers=none,aboveskip=0.5em,belowskip=0.3em]
$( Minimal test: compressed proof with essential hypotheses AND Z marker $)
$( Based on demo0.mm structure - tests the combination that fails in big-unifier.mm $)
$(
  This test case is designed to isolate the bug where:
  - Compressed proofs work (test57 passes)
  - Normal proofs with e-hyps work (demo0.mm passes)
  - BUT compressed proofs with e-hyps + Z marker FAIL (big-unifier.mm fails)

  This file contains both normal and compressed versions of the same proof,
  allowing direct comparison of verification behavior.
$)

$c 0 + = -> ( ) term wff |- $.
$v t r s P Q $.

tt $f term t $.
tr $f term r $.
ts $f term s $.
wp $f wff P $.
wq $f wff Q $.

$( Define term $)
tze $a term 0 $.
tpl $a term ( t + r ) $.

$( Define wff $)
weq $a wff t = r $.
wim $a wff ( P -> Q ) $.

$( Axioms $)
a1 $a |- ( t = r -> ( t = s -> r = s ) ) $.
a2 $a |- ( t + 0 ) = t $.

$( Modus ponens - has essential hypotheses $)
${
  min $e |- P $.
  maj $e |- ( P -> Q ) $.
  mp $a |- Q $.
$}

$( Normal proof of |- t = t - should work $)
th1_normal $p |- t = t $=
  tt tze tpl tt weq tt tt weq tt a2 tt tze tpl tt weq tt tze tpl tt weq tt tt
  weq wim tt a2 tt tze tpl tt tt a1 mp mp $.

$(
  Compressed proof of |- t = t generated by Metamath program.
  This proof:
  1. Uses compressed format
  2. Uses modus ponens (mp) which has essential hypotheses
  3. Uses Z marker to save/recall intermediate results

  Labels:
    A=tt (mandatory f-hyp)
    B=tze, C=tpl, D=weq, E=a2, F=wim, G=a1, H=mp (explicit labels)

  Proof string: ABCZADZAADZAEZJJKFLIAAGHH

  The Z at position 4 saves something, then J (index 9) recalls it.
  This combines all three ingredients that cause big-unifier.mm to fail.
$)
th1_compressed $p |- t = t $=
  ( tze tpl weq a2 wim a1 mp ) ABCZADZAADZAEZJJKFLIAAGHH $.
\end{lstlisting}

\noindent\textbf{test65\_compressed\_assertion\_fhyp\_order.mm} [ACCEPT] \textit{Spec AppB: compressed proofs with scoped f-hyps}\par

\begin{lstlisting}[language=Metamath,basicstyle=\ttfamily\scriptsize,frame=single,framerule=0.3pt,rulecolor=\color{black!20},backgroundcolor=\color{black!2},xleftmargin=1em,framexleftmargin=0.5em,numbers=none,aboveskip=0.5em,belowskip=0.3em]
$( Unit Test 65: Compressed proof with 3-variable axiom $)
$( Tests: Correctly using an axiom with 3 f-hyps in a compressed proof.
   ax-3var has mandatory hyps: wff ph, wff ps, wff ch, |- ph, |- ps
   We prove |- ps by calling ax-3var with substitution ch -> ps
   IMPORTANT: Everything wrapped in outer ${ $} to avoid top-level $e $)
$( Should accept: True $)

$c wff |- $.
$v ph ps ch $.

${
  wp $f wff ph $.
  wq $f wff ps $.
  wc $f wff ch $.

  $( Axiom: from |- ph and |- ps, conclude |- ch $)
  ax.1 $e |- ph $.
  ax.2 $e |- ps $.
  ax-3var $a |- ch $.

  $( Theorem: prove |- ps by using ax-3var with ch -> ps $)
  $( Mandatory hyps of th: wp(A), wq(B), ax.1(C), ax.2(D), th.1(E), th.2(F)
     Note: ax.1 and ax.2 are mandatory because they're used in the proof!
     Label list: (ax-3var) -> G = ax-3var
     Proof: A B B E F G = push wff ph, wff ps, wff ps, |- ph, |- ps, call ax-3var
     ax-3var with {ph->ph, ps->ps, ch->ps} yields |- ps $)
  th.1 $e |- ph $.
  th.2 $e |- ps $.
  th $p |- ps $= ( ax-3var ) ABBEFG $.
$}
\end{lstlisting}

\noindent\textbf{test66\_compressed\_mand\_hyp\_order.mm} [ACCEPT] \textit{Spec AppB: mandatory hypothesis ordering}\par

\begin{lstlisting}[language=Metamath,basicstyle=\ttfamily\scriptsize,frame=single,framerule=0.3pt,rulecolor=\color{black!20},backgroundcolor=\color{black!2},xleftmargin=1em,framexleftmargin=0.5em,numbers=none,aboveskip=0.5em,belowskip=0.3em]
$( Unit Test 66: Compressed proof - mandatory hypothesis ordering $)
$( Tests: Mandatory hyps are f-hyps (in order) + e-hyps (in order) $)
$( Should accept: True $)

$c wff |- $.
$v ph ps $.
wp $f wff ph $.
wq $f wff ps $.

${
  $( Axiom with 2 essential hyps $)
  min.1 $e |- ph $.
  min.2 $e |- ps $.
  ax-min $a |- ph $.
$}

${
  $( Proof using compressed format $)
  $( Mandatory hyps: wp(A), wq(B), th.1(C), th.2(D). Label list: E=ax-min $)
  th.1 $e |- ph $.
  th.2 $e |- ps $.
  th $p |- ph $= ( ax-min ) ABCDE $.
$}
\end{lstlisting}

\noindent\textbf{test67\_toplevel\_essential.mm} [ACCEPT] \textit{Spec allows top-level \textbackslash{}\$e (knife non-compliant)}\par

\begin{lstlisting}[language=Metamath,basicstyle=\ttfamily\scriptsize,frame=single,framerule=0.3pt,rulecolor=\color{black!20},backgroundcolor=\color{black!2},xleftmargin=1em,framexleftmargin=0.5em,numbers=none,aboveskip=0.5em,belowskip=0.3em]
$( Unit Test 67: Top-level $e statements $)
$( Tests: Essential hypotheses at the outermost scope level
   Spec allows this - no restriction in EBNF or spec text.
   metamath-knife incorrectly REJECTS this (non-compliant).
   Official Metamath correctly ACCEPTS this. $)
$( Should accept: True $)

$c wff |- $.
$v ph ps $.
wp $f wff ph $.
wq $f wff ps $.

$( Top-level essential hypothesis - not inside any ${  $} block $)
hyp1 $e |- ph $.

$( This assertion has hyp1 as a mandatory hypothesis $)
ax1 $a |- ps $.
\end{lstlisting}

\noindent\textbf{test\_axmp\_z.mm} [REJECT] \textit{Broken compressed proof - both verifiers reject}\par

\begin{lstlisting}[language=Metamath,basicstyle=\ttfamily\scriptsize,frame=single,framerule=0.3pt,rulecolor=\color{black!20},backgroundcolor=\color{black!2},xleftmargin=1em,framexleftmargin=0.5em,numbers=none,aboveskip=0.5em,belowskip=0.3em]
$( Minimal test: ax-mp with essential hypotheses and Z marker $)

$c wff |- e ( , ) $.
$v x y $.

wx $f wff x $.
wy $f wff y $.

$( Binary connective $)
wi $a wff e ( x , y ) $.

$( Modus ponens with essential hypotheses $)
${
  ax-mp.1 $e |- x $.
  ax-mp.2 $e |- e ( x , y ) $.
  ax-mp $a |- y $.
$}

$( A simple axiom to use $)
ax1 $a |- e ( x , e ( y , x ) ) $.

$(
  Compressed proof using Z:
  Goal: |- e ( x , e ( y , x ) )

  Labels: wx=0, wy=1, wi=2, ax-mp=3, ax1=4

  Proof steps for normal proof:
  wx wy ax1 => |- e ( x , e ( y , x ) )

  As compressed: ( wi ax-mp ax1 ) ABC
    A=0=wx, B=1=wy, C=2=ax1
$)
th1 $p |- e ( x , e ( y , x ) ) $= ( wi ax-mp ax1 ) ABC $.
\end{lstlisting}

\noindent\textbf{test\_compressed\_syl.mm} [REJECT] \textit{Broken compressed proof - both verifiers reject}\par

\begin{lstlisting}[language=Metamath,basicstyle=\ttfamily\scriptsize,frame=single,framerule=0.3pt,rulecolor=\color{black!20},backgroundcolor=\color{black!2},xleftmargin=1em,framexleftmargin=0.5em,numbers=none,aboveskip=0.5em,belowskip=0.3em]
$( Minimal test for compressed proof with multiple floating hypotheses $)

$c term |- |= $.
$v R S T $.

tR $f term R $.
tS $f term S $.
tT $f term T $.

${
  syl.1 $e |- R |= S $.
  syl.2 $e |- S |= T $.
  $( Syllogism axiom $)
  ax-syl $a |- R |= T $.
$}

${
  th.1 $e |- R |= S $.
  th.2 $e |- S |= T $.
  $( Theorem using compressed proof - should push tR, tS, tT before ax-syl $)
  th $p |- R |= T $= ( ax-syl ) BCA $.
$}
\end{lstlisting}

\noindent\textbf{test\_dv\_multi.mm} [REJECT] \textit{Comment says INVALID PROOF - should be rejected}\par

\begin{lstlisting}[language=Metamath,basicstyle=\ttfamily\scriptsize,frame=single,framerule=0.3pt,rulecolor=\color{black!20},backgroundcolor=\color{black!2},xleftmargin=1em,framexleftmargin=0.5em,numbers=none,aboveskip=0.5em,belowskip=0.3em]
$( Test for multi-variable $d declarations - INVALID PROOF (should be rejected) $)

$c wff |- $.
$v x y z $.

wx $f wff x $.
wy $f wff y $.
wz $f wff z $.

${
  $d x y z $.
  ax $a |- x $.

  ${
    th.1 $e |- y $.
    th $p |- z $= wy wz ax th.1 $.
  $}
$}
\end{lstlisting}

\subsection{Large Test Files (Summary Only)}

The following test files exceed 70 lines and are summarized rather than listed verbatim.

\begin{tabular}{llrl}
\toprule
\textbf{File} & \textbf{Verdict} & \textbf{Lines} & \textbf{Spec Rule} \\
\midrule
\texttt{big-unifier-bad1.mm} & REJECT & 100 & 4.2.1: unification/substitution failure \\
\texttt{big-unifier-bad2.mm} & REJECT & 100 & 4.2.1: unification/substitution failure \\
\texttt{big-unifier-bad3.mm} & REJECT & 100 & 4.2.1: unification/substitution failure \\
\texttt{big-unifier.mm} & ACCEPT & 341 & Valid database per spec \\
\texttt{miu.mm} & ACCEPT & 131 & Valid database per spec \\
\texttt{peano-fixed.mm} & ACCEPT & 859 & Valid database per spec \\
\bottomrule
\end{tabular}

